% =====================================================================
% Disertație template — Universitatea din București, Facultatea de Filosofie
% Conformant cu Ghidul de redactare, 31 iulie 2023.
%
% Pandoc usage (Zettlr / CLI):
%   pandoc input.md -o out.pdf \
%     --template=disertatie.tex --top-level-division=section \
%     --pdf-engine=xelatex
%
% Required YAML frontmatter (example):
%   ---
%   title: "From Local to Portable Rigidity"
%   subtitle: "Demonstrative Reference and the Cognitive Ground of Rigid Designation"
%   author: "Florin Cojocariu"
%   program: "Master in Analytic Philosophy"
%   advisor: "Mircea Dumitru"
%   advisor-title: "Prof. Dr."
%   work-type: "Lucrare de disertație"   # optional, this is the default
%   date: "2026"
%   abstract: |
%     Two paragraphs at most.
%   ---
%
% NOTE: The cover page below is HARDCODED (does not read YAML). Edit the
% cover block directly to change title/subtitle/advisor/date on the cover.
% The \title/\author/\date metadata commands are retained but unused on
% the cover.
% =====================================================================

\documentclass[12pt,a4paper]{article}

% —— Core packages ——
\usepackage[utf8]{inputenc}
\usepackage[T1]{fontenc}
\usepackage{textcomp}
\usepackage{times}                       % Times New Roman per guide
\usepackage{amsmath,amssymb}
\usepackage{stmaryrd}
\usepackage[a4paper,left=2.5cm,right=2cm,top=2cm,bottom=2cm]{geometry}
\usepackage{microtype}
\emergencystretch=1em
\usepackage{setspace}
\onehalfspacing                          % 1.5 line spacing per guide
\usepackage[most]{tcolorbox}
\tcbuselibrary{theorems,skins,breakable}

% Optional: Romanian hyphenation/quotes — uncomment if writing in Romanian
% \usepackage[romanian]{babel}

% —— Modal logic, Greek, set-theoretic symbols ——
\DeclareUnicodeCharacter{03B9}{\ensuremath{\iota}}
\DeclareUnicodeCharacter{25A1}{\ensuremath{\Box}}
\DeclareUnicodeCharacter{25C7}{\ensuremath{\Diamond}}
\DeclareUnicodeCharacter{03BB}{\ensuremath{\lambda}}
\DeclareUnicodeCharacter{03C6}{\ensuremath{\varphi}}
\DeclareUnicodeCharacter{00AC}{\ensuremath{\neg}}
\DeclareUnicodeCharacter{22A2}{\ensuremath{\vdash}}
\DeclareUnicodeCharacter{2203}{\ensuremath{\exists}}
\DeclareUnicodeCharacter{2200}{\ensuremath{\forall}}
\DeclareUnicodeCharacter{2227}{\ensuremath{\wedge}}
\DeclareUnicodeCharacter{2228}{\ensuremath{\vee}}
\DeclareUnicodeCharacter{2192}{\ensuremath{\rightarrow}}
\DeclareUnicodeCharacter{2194}{\ensuremath{\leftrightarrow}}
\DeclareUnicodeCharacter{2234}{\ensuremath{\therefore}}
\DeclareUnicodeCharacter{2248}{\ensuremath{\approx}}
\DeclareUnicodeCharacter{2208}{\ensuremath{\in}}
\DeclareUnicodeCharacter{2260}{\ensuremath{\neq}}
\DeclareUnicodeCharacter{2713}{\checkmark}
\DeclareUnicodeCharacter{2212}{-}

% —— Definition box ——
\newenvironment{definitionbox}[1][]{%
  \begin{center}
  \begin{tcolorbox}[
    enhanced, breakable, width=0.8\textwidth,
    colback=gray!2, colframe=gray!60, coltitle=black,
    fonttitle=\bfseries, title={Definition: #1},
    boxed title style={colback=gray!3, colframe=black!60,
      sharp corners, boxrule=0.5pt},
    boxrule=0.5pt, sharp corners,
    left=6pt, right=6pt, top=6pt, bottom=6pt,
    before skip=12pt, after skip=12pt,
  ]}%
  {\end{tcolorbox}\end{center}}

% —— Tables ——
\usepackage{array}
\usepackage{longtable}
\usepackage{booktabs}
\usepackage{calc}
\newcolumntype{C}[1]{>{\raggedright\arraybackslash}p{\dimexpr#1\linewidth-4\tabcolsep\relax}}

% —— Paragraph formatting per guide ——
% First-line indent 1.27 cm on EVERY paragraph (including those that
% follow a heading — indentfirst ensures that).
\usepackage{indentfirst}
\setlength{\parindent}{1.27cm}
\setlength{\parskip}{0pt}
\raggedbottom

% —— Heading formatting per guide ——
\usepackage{titlesec}
\usepackage{needspace}
\usepackage{etoolbox}

\setcounter{secnumdepth}{3}

% \section corresponds to a Chapter (Markdown H1 via --top-level-division=section).
% Guide: bold, centered, each new chapter starts on a new page.
\titleformat{\section}[block]
  {\normalfont\Large\bfseries\filcenter}
  {\thesection.}{0.5em}{}
\titlespacing*{\section}{0pt}{0pt}{1.5em}
\pretocmd{\section}{\clearpage}{}{}

% \subsection corresponds to a Section (H2).
% Guide: bold, indented 1.27 cm from left margin, NOT centered.
\titleformat{\subsection}[hang]
  {\normalfont\normalsize\bfseries}
  {\thesubsection.}{0.5em}{}
\titlespacing*{\subsection}{1.27cm}{1.4em}{0.6em}

% \subsubsection corresponds to a Subsection (H3) — same convention.
\titleformat{\subsubsection}[hang]
  {\normalfont\normalsize\bfseries}
  {\thesubsubsection.}{0.5em}{}
\titlespacing*{\subsubsection}{1.27cm}{1em}{0.4em}

% —— TOC formatting ——
\usepackage{tocloft}
\renewcommand{\cftsecfont}{\small}
\renewcommand{\cftsubsecfont}{\small}
\renewcommand{\cftsubsubsecfont}{\small}
\renewcommand{\cftsecpagefont}{\small}
\renewcommand{\cftsubsecpagefont}{\small}
\renewcommand{\cftsubsubsecpagefont}{\small}
\setlength{\cftbeforesecskip}{0.3em}
\setlength{\cftbeforesubsecskip}{0.1em}
\setlength{\cftbeforesubsubsecskip}{0.1em}

% —— Lists ——
\usepackage{enumitem}
\setlist{nosep, leftmargin=2em, topsep=0.4em, partopsep=0pt,
         parsep=0.2em, itemsep=0.2em}
\setlist[itemize,1]{label=\textbullet}
\setlist[itemize,2]{label=\textendash}

% —— Figures ——
\usepackage{graphicx}
\usepackage{float}
\setkeys{Gin}{keepaspectratio,width=0.7\linewidth}
\makeatletter
\def\fps@figure{H}
\makeatother
\graphicspath{{figures/}}

% —— URLs / DOIs ——
\usepackage{xurl}

% —— CSL bibliography ——
\newlength{\cslhangindent}
\setlength{\cslhangindent}{1.5em}
\newenvironment{CSLReferences}[2]%
  {\begin{enumerate}[label={[\arabic*]},leftmargin=\cslhangindent,itemsep=0pt]%
   \raggedright\sloppy
   \renewcommand{\bibitem}[2][]{\item}}%
  {\end{enumerate}}
\providecommand{\CSLLeftMargin}[1]{\item[#1]}
\providecommand{\CSLRightInline}[1]{#1}

% —— Hyperref (near last) ——
\usepackage[colorlinks=true, linkcolor=black, urlcolor=blue, citecolor=blue]{hyperref}
\urlstyle{same}

% —— Pandoc / longtable compatibility ——
\makeatletter
\@ifundefined{c@none}{\newcounter{none}}{}
\makeatother
\renewcommand{\thenone}{}
\providecommand*{\theHnone}{\arabic{none}}
\providecommand{\tightlist}{}
\providecommand{\pandocbounded}[1]{#1}

% —— Citeproc compatibility ——
\makeatletter
\providecommand{\citeproctext}{}
\let\@oldprotect\protect
\def\@citeprocprotection{\let\protect\@oldprotect}
\newcommand{\citeproc}[2]{#2}
\protected\def\citeprocfootnote#1{\footnote{#1}}
\makeatother

% —— Block quotes ——
\renewenvironment{quote}{%
  \par\vskip 0.5em
  \leftskip=2em\rightskip=2em\small
  \noindent\ignorespaces
}{\par\vskip 0.5em}

% —— Custom notation (PRU) ——
\newcommand{\llangle}{\langle\!\langle}
\newcommand{\rrangle}{\rangle\!\rangle}
\DeclareUnicodeCharacter{27EA}{\ensuremath{\llangle}}
\DeclareUnicodeCharacter{27EB}{\ensuremath{\rrangle}}
\newcommand{\Rfr}[2]{\ensuremath{\mathcal{R}(#1, #2)}}
\newcommand{\LPC}[1]{\ensuremath{\llangle#1\rrangle}}

% —— Document metadata (retained but unused on the hardcoded cover) ——
\title{From Local to Portable Rigidity -\/- Demonstratives and the Metasemantic Theory of Names}
\author{Florin Cojocariu}
\date{}

\begin{document}

% Continuous arabic numbering starting at the cover; display is suppressed
% on front matter pages and resumes at the body (the Introduction).
\pagenumbering{arabic}
\pagestyle{empty}

% =====================================================================
% COVER PAGE (page 1, number hidden) — Anexa II expanded per III.2.a
% =====================================================================
\begin{center}
\vspace*{0.5cm}
{\large UNIVERSITATEA DIN BUCUREȘTI}\\[0.3em]
{\large FACULTATEA DE FILOSOFIE}

\vfill

{\large\bfseries Lucrare de disertație}

\vspace{2em}

{\LARGE\bfseries From Local to Portable Rigidity}\\[1em]
{\Large Demonstrative Reference and the Cognitive Ground of Rigid Designation}

\vfill
\end{center}
\noindent
\begin{minipage}[t]{0.45\textwidth}
\raggedright
\textbf{Coordonator științific:}\\
Prof. Dr. Mircea Dumitru
\end{minipage}%
\hfill
\begin{minipage}[t]{0.45\textwidth}
\raggedleft
\textbf{Absolvent:}\\
Florin Cojocariu
\end{minipage}
\vspace{1.5cm}
\begin{center}
{\large București}\\[0.3em]
{\large 2026-06-15}
\end{center}
\clearpage

% =====================================================================
% ABSTRACT (page hidden)
% =====================================================================

% =====================================================================
% TABLE OF CONTENTS (pages hidden)
% =====================================================================
\setcounter{tocdepth}{3}
\tableofcontents
% tocloft forces \thispagestyle{plain} on the first TOC page; override it.
% (If the TOC ever spans more than one page, also wrap with \pagestyle{empty}
% on each page — handled here because outer \pagestyle{empty} is still active.)
\thispagestyle{empty}
\clearpage

% =====================================================================
% BODY — page numbers become visible from here
% =====================================================================
\pagestyle{plain}

\section*{Abstract}\label{abstract}
\addcontentsline{toc}{section}{Abstract}

The dissertation argues that Kripke's account of rigid designation rests on a cognitive precondition his semantics presupposes but does not describe. The semantic results of \emph{Naming and Necessity} specify what rigid designators do --- a name holds its referent fixed across every world in which that referent exists --- and Kripke argues for those results on grounds that owe nothing to the necessity-of-identity formula. What the results do not specify is the cognitive operation a speaker performs such that her words behave as the semantics describes. Kripke's causal-historical picture of baptism and transmission names this operation from the outside without describing it, and offers itself, in his own words, as a better picture rather than a theory. The dissertation describes the operation.

The proposal starts from a more primitive phenomenon than Kripkean rigidity, which it calls \emph{local rigidity}. When a speaker tracks a particular through encounter and uses a word to pick it out, the word coordinates with that particular directly, without descriptive mediation; within the modal scope of such a use there is nothing for scope to interact with, and the word holds its object fixed across whatever is imagined about it. This is rigidity in the demonstrative case, constituted by the coordination rather than built into the word as a standing property. The dissertation's central claim is that Kripkean rigidity is not primitive but is local rigidity made portable: what a proper name does is carry the singular anchoring first visible in demonstrative encounter into conditions under which the encounter is no longer available.

The argument runs through Russell, Kripke, and the post-Kripkean cognitive literature. Russell identified the relevant operation --- direct, non-descriptive reference under acquaintance --- and then restricted it to incorrigible sense-data on Cartesian grounds, leaving it unavailable for ordinary public names. Prising apart two properties Russell ran together, directness and infallibility, frees the operation for ordinary objects at the cost of nothing the account needs. Among more recent cognitive accounts, Evans's discriminating conception, Recanati's mental files, and Dickie's account of aboutness each capture part of the explanatory burden while leaving the operation itself undescribed. The dissertation locates that operation in object-mode coordination --- a word's grip on an experienced particular --- and treats baptism as the act that stabilizes such a grip into a public name, with testimonial transmission and revisable description carrying the name through the bearer's absence.

The account is metasemantic, not a rival semantics. It does not contest Kripke's model or propose a different one; it explains how a speaker comes to be positioned such that the model's stipulation correctly describes her practice. Its scope is the single speaker's anchoring of a word to a particular --- the question of what one speaker does in getting a grip on an individual --- and not the distinct question of how a community of speakers coordinates on a shared referent, which the dissertation marks as a connected problem and reserves for separate work. Within that scope, the result is narrow and specific: not a theory of reference, but the closing of one gap in an existing theory, by describing the cognitive operation that Kripke's semantics presupposes and Russell located before placing it out of reach.

\section{Introduction}\label{introduction}

Saul Kripke's arguments in \emph{Naming and Necessity} and \emph{Identity and Necessity} establish that proper names are rigid designators: a name picks out the same individual in every possible world in which that individual exists. The arguments are compelling and the conclusion is by now a standard starting point. What the arguments establish is what rigid designators do. What they do not establish, and what Kripke does not set out to establish, is the cognitive operation that makes the doing possible.

Reference to a particular is not generated by linguistic machinery operating on itself. No tightening of the network of words produces a grip on an individual; descriptions locate whatever satisfies a condition, not an individual as such. The grip comes from another direction --- from using a word where it meets an encounter. Kripke saw this, which is why his positive account appeals not to descriptions but to an \emph{initial baptism} and a \emph{chain of communication}. A baptismal act ties a word to an individual; a chain of communication preserves the tie. Neither is a logical relation. Kripke is clear about this: both are causal-social practices, and his account is a description of them. What his account does not describe is the cognitive operation those practices implement --- what a speaker does, at the baptism and afterward, such that the name holds to its bearer. This dissertation describes that operation. The aim is not to correct Kripke but to extend the explanatory target, from the semantic and social profile of rigid designation to its cognitive realization.

The thesis is that Kripkean rigidity is not primitive. It rests on a more basic phenomenon, present whenever a speaker tracks a particular through encounter and uses a word to pick it out: the word coordinates with that particular directly, and holds it fixed across whatever is imagined about it. Call this \emph{local rigidity}. A proper name is what carries the singular anchoring first visible in such an encounter into conditions where the encounter is no longer available --- where the bearer is absent, the speaker is not the one who named, and the chain of use is long. Baptism is the act that stabilizes a live coordination into a public name; transmission and revisable description carry the name across the absence. Kripkean rigidity, on this account, is local rigidity made portable.

Two questions must be kept apart, because the dissertation answers only one. The first is how a single speaker gets a grip on a particular at all --- the operation underneath one rigid use, the question pursued here. The second is how a community of speakers stays coordinated on a shared referent across time --- how reference is negotiated and sustained between users. The second question is real and has a substantial literature, but it presupposes the first: there must be a grip on a particular before there is anything for a community to coordinate on. The dissertation isolates the first, the in-the-moment question, and treats the in-time problem of inter-speaker coordination as a connected matter reserved for separate work. Its scope is singular reference; general-term rigidity, fictional and empty names, and mathematical reference are marked where they arise as open questions rather than treated.

The argument is sequential rather than a battery of independent proofs. It begins by laying out the field of direct-reference theory and the rigid-designation thesis in their own terms, then locates the cognitive gap at a specific and resistible point in Kripke's own argument, then shows that the most developed existing attempts to fill the gap each leave the central operation undescribed. It then introduces the vocabulary in which the operation can be described, gives the constructive account --- local rigidity, baptismal stabilization, testimonial transmission, attribution across absence --- and defends it against the objections it most invites. The account is throughout metasemantic, not a rival semantics: it leaves Kripke's formal results untouched and explains how a speaker comes to be positioned such that those results correctly describe her practice.

\section{Direct Reference and the Field}\label{direct-reference-and-the-field}

The thesis that a name contributes its bearer, and not a descriptive condition, to what is said is older than the apparatus of rigid designation and is in principle separable from it. It is worth stating on its own, because the dissertation's question is pitched at the level of this thesis rather than at the level of the modal semantics that later regiments it.

The thesis is direct reference: the semantic contribution of certain expressions is their referent itself, with no Fregean sense interposed between word and object. Demonstratives and indexicals are the paradigm --- terms which, as Kaplan puts it, refer directly without the mediation of a Fregean sense as meaning. The point is not that nothing whatever connects the term to its object; a directly referential term refers relative to a context of utterance, and context plainly does mediation of some kind. The point is that what the term contributes to the proposition is the object, not a condition the object satisfies. The position took shape as a rejection of the descriptivist inheritance from Russell, on which a name goes proxy for a cluster of identifying conditions, and it reached its sharpest form in Kaplan's separation of two kinds of meaning: a \emph{character} fixed by linguistic rule and a \emph{content} that, for a directly referential term, is simply the referent.

The line to this thesis runs through Donnellan, and Donnellan's two contributions are distinct. The first is the distinction between two uses of a definite description. Used \emph{attributively}, a description states something about whoever or whatever is the so-and-so; used \emph{referentially}, it serves to enable the audience to pick out whom the speaker has in mind, and states something about that person or thing. In the attributive use the description occurs essentially: the claim is about whatever fits it. In the referential use the description is a tool for directing attention to an individual the speaker already has in view, and reference can succeed even where the object does not satisfy the words --- the man taken to be Smith's murderer may be innocent, yet the speaker has said something about him. The moral Donnellan draws is that expressions denote, while people refer: denoting is a relation a form of words bears to whatever fits it; referring is something a speaker does in using words to talk about a particular.

Donnellan's second contribution concerns proper names directly, and here it is easy to over-assimilate him to Kripke. Both reject descriptivism and both make the reference of a name turn on history rather than on the descriptions a speaker can supply. But the histories they invoke differ. Kripke emphasizes a name-token's causal connection to previous tokens of the name; Donnellan focuses instead on the historical-causal explanation of the speaker's present use and the mental state that produces it. Donnellan offered this as a historical-explanation account and was candid that he had no general theory of which historical connection, among the many linking an utterance to things, is the reference-fixing one. The gap the dissertation addresses is visible already here, as the place where a historical connection is named but its cognitive operation is left unspecified.

These are the materials of the field the dissertation enters. Direct reference fixes the explanatory target --- reference to a particular, not satisfaction of a condition. Donnellan supplies the datum that reference is something speakers do with particulars they have in mind, and the further datum that names trace to their bearers through a history neither speaker nor theory has fully described. Kripke's rigid-designation thesis, to which the next chapter turns, is the semantic regimentation of this same target. A reader may hold that the field is exhausted by direct-reference and coordination theory, and that once reference is traced to a particular and sustained across a community no further question remains. The dissertation's claim is that one question does remain, and it is prior to both: what none of these accounts isolates is the cognitive operation by which a speaker gets the particular in mind in the first place, such that a name can later be anchored to it. That operation is the dissertation's subject.

One adjacent body of work has to be set aside deliberately rather than passed over, because it addresses a real problem the present account does not take up---one that arises naturally once you try to supply a concrete mechanism for the causal-historical chains that Donnellan and Kripke invoke in different ways. Beginning with Evans and continued by Fine, a strand of the literature takes the central phenomenon of reference to be coordination between language-users. Evans's causal theory of names already locates the reference of a name not in any single baptism but in the dominant source feeding a name-using practice --- which is why a name can drift, as \emph{Madagascar} did, from the mainland region it first marked to the island it now marks, as the practice's centre of gravity shifts. Fine takes this further, making coordination the primary semantic relation: a relation of representing-as-the-same that holds between expressions and does not reduce to their intrinsic features. On Fine's account of names, coordination is inter-subjective --- when one speaker derives her use of a name from another's, the tokens are tied together through a common communal use.

These accounts describe something real: how speakers stay aligned on a referent across a community, how a name's reference is negotiated and sustained between users, how co-reference becomes a semantic fact in its own right. But this is a different question from the one pursued here. The coordination literature runs between speakers, in time: it asks how many users converge on one referent. The present account runs in the moment, between a single speaker's use and the particular it is anchored to: it asks what one speaker does in getting a grip on an individual at all. The in-time question presupposes the in-the-moment one --- there must be a grip on a particular before there is anything for a community to coordinate on---and it is the in-the-moment question, the cognitive operation underneath a single rigid use, that the dissertation isolates. The in-time problem of inter-speaker coordination, including whether Fine's relationism is the right treatment of it, is left to separate work.

\section{Rigid Designation}\label{rigid-designation}

This chapter sets out Kripke's thesis that proper names are rigid designators and the arguments he gives for it, in his own terms. It presents those arguments, then lays out why descriptions, taken as descriptions, do not supply the kind of rigidity the arguments demand. The arguments establish a semantic property of names; what they do not yet supply is an account of the cognitive operation by which speakers come to use a name in the way the thesis attributes to it. Marking that the arguments leave this open is the work of the next chapter; supplying it is the work of the chapters after.

\subsection{Kripke's Independent Argument}\label{kripkes-independent-argument}

Kripke's thesis is that proper names are rigid designators: a proper name designates the same individual in every possible world in which that individual exists. \emph{Naming and Necessity} defends the thesis through three converging arguments, none of which depends on the necessity-of-identity formula. The formula regiments the modal consequence once rigidity is granted; the case for rigidity runs separately.

The modal argument turns on counterfactual intuition. Aristotle could have died in infancy, never taught Alexander, and never written the \emph{Nicomachean Ethics}. In that counterfactual it is Aristotle who failed to do these things, not someone else. If ``Aristotle'' abbreviated ``the teacher of Alexander who wrote the \emph{Nicomachean Ethics},'' then in a world where no one satisfied that description the name would refer to no one; if it referred to whoever did satisfy it, the counterfactual would be about someone else. Both readings clash with the intuition that the counterfactual is about the man Aristotle himself.

The epistemic argument turns on ignorance and error. A competent user of ``Feynman'' may know only that Feynman was a physicist, or may believe things about him that are false. On a descriptivist theory, such a user either fails to refer or refers to whoever fits her associated description. The sharpest version is the Gödel/Schmidt case: suppose the incompleteness proof we attribute to Gödel was in fact discovered by a man named Schmidt, and Gödel merely stole it. On a descriptivist theory built around ``the prover of the incompleteness of arithmetic,'' ``Gödel'' would refer to Schmidt. The intuition is that ``Gödel'' still refers to Gödel. Reference runs through the name, not through the descriptions attached to it.

The semantic argument targets substitutivity in modal contexts. ``Necessarily, Benjamin Franklin is Benjamin Franklin'' is true; ``Necessarily, the inventor of bifocals is the inventor of bifocals'' is also true; but ``Necessarily, Benjamin Franklin is the inventor of bifocals'' is false --- he might never have invented them. The name and the description share an actual-world referent but differ in modal profile. They cannot be synonyms.

Taken together, the three arguments support the distinction between \emph{de jure} rigidity, characteristic of proper names and rigid by how the designator functions, and \emph{de facto} rigidity, which can attach to descriptions whose satisfiers happen to be necessarily fixed across worlds, as with ``the least prime.''\footnote{Kripke, \emph{Naming and Necessity}, p.~21, fn. 21. The \emph{de jure} / \emph{de facto} distinction matters here because the formula's presupposition is \emph{de jure} rigidity --- rigidity by how the designator functions, not by metaphysical accident.} What they do not by themselves supply is an account of the cognitive operation by which speakers come to use a name in the rigid way the semantic thesis attributes to it.

\subsection{Why Descriptions Cannot Supply Rigidity}\label{why-descriptions-cannot-supply-rigidity}

Descriptions, qua descriptions, pick out satisfiers world by world. ``The inventor of bifocals'' picks out whoever, in a given world, invented bifocals; different worlds can therefore yield different referents. This world-relative behavior is constitutive of descriptive terms: they specify a condition, and reference goes wherever the condition is met. Adding further descriptive material may narrow the range of worlds in which the description is satisfied, but it does not by itself make the reference invariant across those worlds.

The formal expression of the difficulty is the scope problem. In \(\Box \, F(\iota x\, G x)\), the description \(\iota x\, G x\) can take scope inside or outside the modal operator. Narrow scope (``necessarily, \(F\) of whoever is \(G\)'') requires that in every world the satisfier of \(G\) in that world is also \(F\). Wide scope (``of whoever is actually \(G\), necessarily \(F\)'') requires that the actual-world satisfier of \(G\) is \(F\) in every world. The wide-scope reading behaves modally as if the description were rigid: the referent is fixed in the actual world and held constant. But the rigidity is positional, not constitutive --- it is achieved by choice of scope, not by anything intrinsic to the description.\footnote{Kripke, \emph{Naming and Necessity}, p.~12, fn. 15. Kripke notes that taking a description outside the scope of a modal operator can produce something that behaves like rigidity, and is explicit that this is not what the rigid-designation thesis is asserting.}

Something similar applies to Kaplan's \(\text{dthat}\) operator. \(\text{dthat}(\iota x\, F x)\) takes a description and yields a directly referring term whose reference is fixed in the actual world and then held constant across worlds. The result behaves rigidly, but the rigidity is introduced by the rigidifying device, not supplied by the descriptive content itself.\footnote{A family of two-dimensional treatments refines the classification of rigidity beyond the de jure/de facto and scope distinctions --- Davies and Humberstone's fixedly operator, Kaplan's character/content separation, Chalmers's primary and secondary intensions. These sharpen what kinds of world-invariance a designator can have; none addresses the cognitive operation by which a speaker uses a name as rigid, which is the question the dissertation takes up.}

Names differ from rigidified descriptions in a further respect. A proper name does not contain descriptive structure that could generate a genuine narrow-scope reading of the ordinary descriptivist kind. ``Aristotle'' does not pick out whoever happens to satisfy some condition in each world; if it is rigid, it is rigid by the way it functions, not by scope placement or post hoc rigidification.

\subsection{Where This Leaves the Argument}\label{where-this-leaves-the-argument}

The arguments of this chapter establish a semantic property: names hold their referents fixed across counterfactual variation, and descriptions, even when scope or a rigidifying device makes them behave rigidly, do not have that property constitutively. What the arguments characterize is what rigid designators do. What they do not characterize is the cognitive operation by which a speaker comes to use a name in this way --- what a speaker does, at baptism and afterward, such that the name holds its referent fixed. A reader may accept everything argued here and hold that no such question arises: that the semantics is the whole story, and that asking what speakers do to satisfy it is a confusion of levels. The next chapter argues that the question is real, and that it is left open at a specific and locatable point in Kripke's own argument.

\section{The Problem}\label{the-problem}

The previous chapter argued for the rigidity of names through three intuitive arguments, none of which turns on the necessity-of-identity formula. There is a further argument, the one most often associated with Kripke's result, that runs through that formula directly. Looking closely at it is worthwhile not because it adds to the case for rigidity --- it does not --- but because it is where the cognitive question this dissertation pursues becomes visible. The formula is valid; the philosophical thesis it is used to support is about names; and the move between the two passes through a place where a cognitive mechanism is presupposed and not supplied. This chapter locates that place.

Kripke's necessity-of-identity argument has two parts: a formal core, the modal logic of identity under fixed assignment, \(\forall x \forall y (x = y \to \Box(x = y))\); and a metasemantic premise, that ordinary proper names function as rigid designators. The formal core delivers the modal consequence cleanly. The metasemantic premise carries the philosophical weight, and Kripke argues for it independently of the formula, on intuitive grounds --- the counterfactual cases of \emph{Naming and Necessity} and the Gödel/Schmidt thought experiment.\footnote{Kripke, \emph{Naming and Necessity} (1980), pp.~100--105; ``Identity and Necessity'' (1971). The reading of the formula as a conditional whose force depends on the rigidity of its terms is standard; see LaPorte, ``Rigid Designators,'' \emph{Stanford Encyclopedia of Philosophy}, §1.1. The intuitive cases for the rigidity of names: \emph{Naming and Necessity}, pp.~74--78 (Nixon), pp.~60--62 and 75--76 (Aristotle), pp.~83--87 (Gödel/Schmidt).}

This split is a standard reading in the literature, and Kripke himself frames the argument in these terms. The chapter takes this as its starting point and asks a finer question: what does the metasemantic premise commit us to, once we look closely at \emph{how the formula is applied to natural-language cases?} The argument proceeds in three steps. The first specifies what the formal core does and does not establish. The second draws a distinction within the metasemantic premise itself: the word ``object,'' used in Kripke's own gloss and throughout the secondary literature, runs together three different kinds of things, and the formula's validity holds for one while the philosophical conclusion is about a third. The third locates where a cognitive mechanism is needed if the metasemantic premise is to be more than a posit.

The chapter does not supply the mechanism. That is the work of the constructive chapters. What it does is mark the place where the mechanism is needed, with enough precision that the proposal can later be evaluated against the right desideratum.

\subsection{What the Formula Actually Does}\label{what-the-formula-actually-does}

The formal core of the argument is uncontroversial. Under any assignment function \(g\) in a Kripke model, if \(g(x) = g(y)\), then in every world the model constructs, \(g(x) = g(y)\). The assignment function does not vary across worlds in the standard semantics; once an assignment is fixed, the variables continue to pick out their assigned values across every world the operator \(\Box\) quantifies over. The formula is valid by construction.

Two observations follow. They are unremarkable in themselves; together they specify what work the formula can and cannot do for the philosophical project.

The formula's validity rests on the rigidity of variables under assignment, which is a feature of the formal apparatus, not a feature of natural-language names. When Kripke applies the formula to names, he is applying a formal result whose validity holds for one kind of designator (variables under assignment) to another kind of designator (proper names). The application requires the assumption that names exhibit the same world-invariance variables have by construction. \emph{This is the rigid-designation thesis itself, not a consequence of it.} Kripke himself distinguishes the rigidity of names --- \emph{de jure} rigidity, built into the kind of designator a name is --- from \emph{de facto} rigidity, which attaches to descriptions whose satisfiers happen to be necessarily fixed (such as ``the least prime'').\footnote{Kripke, \emph{Naming and Necessity}, p.~21, fn. 21. The \emph{de jure} / \emph{de facto} distinction matters here because the formula's presupposition is \emph{de jure} rigidity --- rigidity by how the designator functions, not by metaphysical accident.} He is also explicit that the wide-scope behavior of descriptions in modal contexts is not what the rigid-designation thesis is asserting: the thesis is about the kind of designator names are, not about scope manipulation.\footnote{Kripke, \emph{Naming and Necessity}, p.~12, fn. 15. Kripke notes that taking a description outside the scope of a modal operator can produce something that behaves like rigidity, and is explicit that this is not what the rigid-designation thesis is asserting.}

The formula has no resources for the intensional dimension of the cases that motivate Kripke's project. Hesperus/Phosphorus, Benjamin Franklin and the inventor of bifocals, water and H\(_2\)O --- these are cases where identity is empirically discovered and then recognized as necessary. The formula regiments the modal status once rigidity is granted; the epistemic route by which speakers reach the identity, and the cognitive capacities involved in using a name as rigid in the first place, lie outside what the formula can address. The formula is extensionally valid but intensionally silent.\footnote{A constraint from the \emph{Tractatus} bears on what kind of identity claim is at stake. Wittgenstein observes (5.5303) that ``to say of two things that they are identical is nonsense, and to say of one thing that it is identical with itself is to say nothing at all.'' \emph{Identity claims are claims about designators co-referring, not about objects themselves being one.} This is why the necessity-of-identity formula, applied to names, concerns the modal status of co-reference between designators rather than a metaphysical relation among objects --- and it is co-reference between distinct designators that the empirical cases (``Hesperus is Phosphorus'') turn on. The bearing of 5.5303 on Kripke's necessity-of-identity argument specifically~is less often discussed than~the rule-following debate through which Wittgenstein and Kripke are usually connected, though I do not claim it has gone unnoticed.}

\subsection{Three Levels of ``Object''}\label{three-levels-of-object}

The standard reading distinguishes the formal core from the metasemantic premise. The chapter adds a finer distinction within the metasemantic premise itself. When the formula is applied to ordinary identity claims, the word ``object'' --- used in Kripke's own gloss and throughout the secondary literature --- covers three distinct things, and the philosophical work happens at a different level than the formal validity does.

\emph{Domain elements} are mathematical objects: points in a set, self-identical by set-theoretic construction, their identity invariant across any world the model constructs. The formula's validity operates at this level. Variables are assigned to domain elements; the assignment function is the device through which formal rigidity is achieved.

\emph{Encountered particulars} are individuals as they figure in perception and re-identification: this dog here, that star there, the man one might have met. They are individuated through encounter; their identity across encounters is an epistemic achievement, with room for error. The early astronomer saw Hesperus in the evening sky and Phosphorus in the morning sky; the identity of the two had to be established, not given with the encounters. Encountered particulars are the bridge between the formal apparatus and the philosophical applications, because they have the concreteness that descriptions can apply to (so the cases Kripke discusses are intelligible at all) while also being the sort of thing that can bear a name.

\emph{Named individuals} are bearers of proper names: Aristotle, Venus, Hesperus, Benjamin Franklin. They are individuated through baptism and tracked across absence through testimonial networks. Their identity across different names is discoverable --- the Hesperus/Phosphorus case is the paradigm. The philosophical thesis of rigid designation operates at this level: names hold their referents fixed across counterfactual variation, where descriptions could track different individuals in different worlds.

Kripke is of course aware that the formal result concerns variables under assignment, not names as such. The point, rather, is that the philosophical application requires a bridging premise: ordinary names must function with the relevant kind of referential stability if the formal result is to illuminate natural-language identity statements. My claim is therefore not that Kripke commits a simple slide, but that the cognitive realization of the bridge remains undertheorized.

What the passage does is move smoothly between these levels in English. The smoothness can obscure distinct explanatory levels. Each level handles part of the argument; no single level handles the whole. The metasemantic premise --- that ordinary names exhibit the world-invariance variables have by construction --- is what makes the move from the first level to the third coherent. The premise is supplied by the rigid-designation thesis, which Kripke argues for separately on intuitive grounds (the modal-intuition tests of \emph{Naming and Necessity}, the counterfactual cases about Nixon and Aristotle).\footnote{A finer-grained version of this analysis distinguishes four readings of ``object x'' in the passage --- mathematical, concrete, positional, name-substitute. That typology refines the three-level distinction without altering its upshot, and is not needed for the argument here.}

The three-level distinction is useful for the rest of the dissertation. The cognitive question --- what enables speakers to use names as rigid designators --- is a question about how the metasemantic premise gets satisfied at the level of named individuals. The mechanism proposed later in the dissertation operates between encountered particulars and named individuals: it explains how a baptism can establish a name whose later use tracks the same referent across absence. The formal apparatus does not need a cognitive story; the philosophical application does.

\subsection{The Cognitive-Mechanism Gap}\label{the-cognitive-mechanism-gap}

What the standard reading leaves open is the cognitive question. Kripke argues that names are rigid; he does not argue that they are rigid because of any particular cognitive mechanism. The causal-historical sketch in \emph{Naming and Necessity} --- baptism fixes the reference, the name propagates through chains of communication --- is offered explicitly as a ``better picture'' rather than as a theory. Kripke is unusually direct about this self-positioning: he is not, he says, ``trying to give any sort of necessary and sufficient conditions for reference\ldots{} in any rigorous sense''; he is ``painting a general picture.''\footnote{Kripke, \emph{Naming and Necessity}, pp.~93--94. The full passage: ``I'm not, in fact, trying to give any sort of necessary and sufficient conditions for reference\ldots{} in any rigorous sense'' (p.~94); ``I think it is a much better picture than the picture presented by a description'' (p.~93).} Kripke gives a causal-social picture of reference-fixing and transmission --- baptism, communicative chains, the intention to preserve reference. What he does not aim to supply is a theory of the cognitive operations that implement these practices in actual speakers. The present account accepts that picture and asks what cognitive coordination it presupposes. This is not a criticism of Kripke's semantic theory; it is an extension of its explanatory target, from semantic and causal-social structure to cognitive implementation.

That a cognitive mechanism is needed at all, and that the available candidates may be too narrow, is something Kripke himself notices in passing. In ``Identity and Necessity,'' he quotes Russell at length on demonstratives as the only true names --- terms that ``really just name an object without describing it'' --- and on acquaintance as the only ground secure enough for identity claims to be ``exempt from all imaginable doubt.'' Kripke's gloss is that on Russell's view, only names of immediate sense data, expressed by demonstratives like ``this'' and ``that,'' qualify as genuine names.\footnote{Kripke, ``Identity and Necessity'' (1971), p.~144. The Russell passage Kripke quotes is from \emph{The Philosophy of Logical Atomism} (1918) and runs: ``if we want to reserve the term `name' for things which really just name an object without describing it, the only real proper names we can have are names of our own immediate sense data, objects of our own immediate acquaintance. The only such names which occur in language are demonstratives like `this' and `that.'\,''} The passage is not used by Kripke to endorse Russell's account; it is used to motivate the contrast with ordinary proper names, which are rigid in Kripke's sense but do not meet Russell's Cartesian standard. What the passage acknowledges, without developing, is that there is a candidate cognitive mechanism for rigidity --- demonstrative grounding under acquaintance --- and that Russell himself made it too narrow to cover the cases Kripke wants to cover. The mechanism Russell identified is the right kind of thing; the scope at which Russell allowed it is too restricted. Whether that mechanism can be preserved while its scope is widened is taken up in the next chapter.

This gap is widely acknowledged in the secondary literature. As the \emph{Stanford Encyclopedia's} discussion of rigid designation notes, Kripke's appeal to stipulation in the treatment of \emph{de jure} rigidity is ``more like a promissory note than the satisfaction of an explanatory obligation.''\footnote{LaPorte, ``Rigid Designators,'' §2.1.1.} The promissory note is for an account of the mechanism by which names achieve, in actual linguistic and cognitive practice, what variables achieve by construction. The note need not be read as three separate debts. Rigid tracking, baptismal fixing, and chain-internal transmission may be understood as three roles that depend on a common capacity: to single out a particular and continue to track it through change, across modal variation, and across handoffs between speakers. The semantic theory describes what this capacity delivers; it does not describe the capacity itself.

\subsection{The Problem Stated}\label{the-problem-stated}

The problem can be put as three claims, set out so that a reader who wishes to resist it can see exactly where to do so.

First, the formal semantics of rigid designation models a referential achievement: it represents a name as holding its referent fixed across every world in which the referent exists. The model is correct about what it models; the dispute is not with the semantics.

Second, the model represents this achievement without supplying the operation that produces it. Holding a referent fixed is something a speaker does --- at the baptism that introduces the name, in retaining the name across the bearer's absence, in inheriting the name from others who introduced it before her. The semantics specifies the output of that doing and is silent on the doing.

Third --- and here the argument turns on a reading of Kripke, not on anything he says outright --- his account captures the output of a cognitive operation without describing the operation. The causal-historical picture is the nearest he comes, and it describes the social shape of the operation from the outside --- who baptizes, who transmits --- not the capacity that each speaker exercises from the inside.

Each claim can be rejected on its own. One could deny the first by holding the semantics itself defective; this dissertation does not. One could deny the second by holding that there is no operation to describe --- that holding-fixed is a feature of the formal apparatus and nothing a speaker does. One could grant the first two and deny the third by holding that the cognitive operation, while real, is not the business of a theory of reference. The dissertation's wager is that the second and third are where a reader should feel the pull the other way: that there is an operation, that it is describable, and that describing it illuminates the semantics rather than competing with it. The chapters that follow describe it.

It is worth saying plainly, before the work of describing the operation begins, what the dissertation does and does not claim to be new. Direct reference is not new, nor is the rigidity of demonstratives, nor the causal-historical picture of how names are transmitted; these are the materials the account works with, not its contribution. What is offered as new is a single claim about how these materials stand to one another: that Kripkean rigidity is not a primitive feature of names but the social portability of a more basic holding-fixed, the one a demonstrative coordination has when it grips a particular in an encounter. The claim is a proposal about the cognitive operation Kripkean practice appears to presuppose, not a demonstration that the operation must be as described; it earns its place, if it does, by organizing the data the rival pictures leave scattered and by predicting where reference survives and fails. Everything that follows is the development and defense of that one derivation claim.

\section{Existing Attempts Fall Short}\label{existing-attempts-fall-short}

The previous chapter located the gap: the semantics of rigid designation describes a referential achievement without describing the operation that produces it. The gap is not unnoticed in the literature, and several accounts have tried to fill it. This chapter examines the most developed of them. The aim is not to refute these accounts --- each identifies something the eventual proposal will need --- but to show that each locates part of the burden while leaving the central capacity undescribed. The chapter takes Russell first, because Russell came closest: he identified the right kind of operation and then, on epistemological grounds, restricted it out of reach. The cognitive accounts that follow --- Evans, Recanati, Dickie --- each capture a different part of what is needed without supplying the whole.

\subsection{Russell's Mechanism and His Retreat}\label{russells-mechanism-and-his-retreat}

If descriptions cannot supply rigidity, the question becomes what non-descriptive cognitive operation can. In a passage from \emph{The Philosophy of Logical Atomism} that Kripke quotes, Russell identifies the operation but immediately restricts its scope:

\begin{quote}
``If we want to reserve the term `name' for things which really just name an object without describing it, the only real proper names we can have are names of our own immediate sense data, objects of our own immediate acquaintance. The only such names which occur in language are demonstratives like `this' and `that.'\,''
\end{quote}

Two claims are joined in the passage, and they must be separated.

The first claim is the isolation of a kind of non-descriptive reference. Russell singles out the demonstrative and characterizes it by what it does not do: it does not refer by way of a descriptive intermediary. This isolates the kind of operation a non-descriptive account of rigidity needs: direct, non-descriptive coordination with a particular.

The second is the retreat. Having identified the mechanism, Russell restricts it to immediate sense data --- the subset of cases in which the reference is exempt from Cartesian doubt. The restriction is epistemological, driven by a foundationalist demand for certainty. Under Russell's standard, even ``this chair'' is not a genuine name, because one could be hallucinating the chair.

The dialectical position of this dissertation begins here. Russell's demonstrative model identifies the kind of operation the rigid-designation literature seems to require: a way for a term to hold its referent fixed without reference being routed through identifying descriptive content. But by confining that operation to incorrigible cases of acquaintance, Russell leaves it unavailable for ordinary public proper names.

The restriction looks unnecessary once two properties are separated: directness and infallibility. For present purposes, a cognitive operation counts as direct when reference is not fixed by the mediation of identifying descriptive content; nothing in that notion by itself entails immunity to error. A direct operation may still misfire under sensory noise, contextual shift, or attentional lapse. Directness is thus a structural feature of the operation, whereas infallibility is an epistemic standard it may or may not satisfy. Once the two are prised apart, Russell's mechanism can be retained without his Cartesian restriction. What Russell saw, and then placed out of reach, is the operation the rest of the dissertation develops --- once it is freed from the demand for certainty that drove him to sense-data.

A reader may resist the separation. One could hold that Russell was right to tie the two together --- that an operation which can misfire is not, in the relevant sense, direct at all, and that dropping infallibility just is dropping the directness Russell isolated. The dissertation's wager is the opposite: that directness is a fact about how reference is routed and infallibility a fact about whether it can err, and that the two come apart cleanly. The argument for the separation is made when the operation is described; here it is only marked as the point on which Russell's retreat turns.

\subsection{The Cognitive Accounts}\label{the-cognitive-accounts}

Russell's is not the only attempt. A body of more recent work takes up the cognitive side of reference directly, and three accounts are the most developed: Evans's discriminating conception, Recanati's mental files, and Dickie's account of aboutness. Each correctly holds that reference to a particular requires a non-descriptive cognitive relation to it, and each describes part of what such a relation involves. None, the chapter argues, isolates the capacity at issue --- the single operation that holds a particular fixed and scales from a live encounter to a name used in the bearer's absence.

Evans requires, for thought about an object, a \emph{discriminating conception}: a non-trivial capacity to single the object out, to tell it apart from all others. This is the right kind of demand --- it locates something on the thinker's side that descriptions alone do not supply, a capacity to discriminate rather than merely a body of information. What it does not do is say what that capacity consists in at the level where it operates. The discriminating conception is stated as a requirement reference must meet, not as the operation by which the requirement is met; it names the achievement and leaves its mechanism open.

Recanati locates singular reference in \emph{mental files}: non-descriptive vehicles of singular thought, individuated not by the information they contain but by the relation the file exploits --- what Recanati calls an epistemically rewarding relation, with perceptual acquaintance as the paradigm. The file refers to whatever stands at the far end of that relation, and the information in it may be false without changing what it is about. This correctly separates reference from descriptive content and ties it to a relation rather than a satisfied condition. But the file is a posited container, and the epistemically rewarding relation is named rather than analyzed: the account tells us that reference runs through such a relation without describing the operation that constitutes it. What does the holding-fixed is attributed to the file; what the file does is left as the explanandum.

Dickie grounds reference in \emph{cognitive focus}: a body of thought is about an object when its means of justification converge on that object, so that the thinker would be unlucky, given how the beliefs are justified, if they failed to match it. Perceptual attention, understanding of a name, and grasp of a description are distinct routes by which this convergence is secured. This is a genuine advance in saying what fixes aboutness --- it gives a criterion, and one that handles perceptual and testimonial cases together. But it pitches the account at the level of justification: it explains in virtue of what a thought is about an object, not what cognitive operation a speaker performs in holding that object fixed across the modal and temporal variation that rigid designation involves. The convergence of justification is a condition on aboutness; it is not yet the tracking operation the gap calls for.

\subsection{What the Survey Shows}\label{what-the-survey-shows}

The accounts do not fail in the same way, and the differences are instructive. Russell found the operation and restricted it to a domain too narrow to cover ordinary names. Evans stated the requirement the operation must meet without giving the operation. Recanati tied reference to a relation and housed it in a file, naming the relation rather than describing what it does. Dickie gave a criterion for aboutness at the level of justification, one level up from the tracking the gap concerns. What none supplies is a single capacity to single out a particular and continue to track it through change, across modal variation, and across the handoff from one speaker to another --- the capacity the semantic theory presupposes and does not describe.

A reader may judge that one of these accounts already discharges the burden --- that Evans's discriminating conception, or Recanati's relation, or Dickie's convergence, is the cognitive story, and that what follows adds only vocabulary. The chapters ahead are the argument that it does not: that each names or constrains the capacity without describing it, and that describing it --- as one operation, present in the demonstrative encounter and carried by the name past it --- does work these accounts leave undone. The next chapter introduces the terms in which that operation is described.

\section{The Functional-Mode Distinction}\label{the-functional-mode-distinction}

The present chapter introduces the analytical vocabulary in terms of which the rest of the dissertation closes the gap. The framework's basic primitives --- the experiential pattern \(\{X\}\), the linguistic pattern \(\langle x \rangle\), and the coordination operator \(\mathcal{R}\) --- are set out in detail in the PRU framework document and are not duplicated here. The chapter develops what they support: a functional-mode distinction between \(x^o\) and \(x^c\), a treatment of predication as coordination rather than function-application, the notion of local rigidity, and re-readings of two classical modal puzzles --- the \emph{de re}/\emph{de dicto} distinction and the necessary \emph{a posteriori}.

\subsection{\texorpdfstring{The \(x^o/x^c\) Functional-Mode Distinction}{The x\^{}o/x\^{}c Functional-Mode Distinction}}\label{the-xoxc-functional-mode-distinction}

A word, on the account developed here, has two functional modes. As \(x^o\), it coordinates with a particular instance of an experiential pattern \(\{X\}\). As \(x^c\), it operates within the linguistic pattern \(\langle x \rangle\) --- the network of utterances, frames, and inferential relations in which the word figures. The distinction is functional rather than lexical: the same word-form, ``bird,'' can play either role, and which role it plays in a given utterance depends on what work the word is being asked to do.

Two examples isolate the modes. In ``that bird is injured,'' ``bird'' functions as \(bird^o\). It coordinates with one bird in the speaker's perceptual field. The utterance is about that bird; substitute another bird and the truth-value can change. In ``birds fly,'' the word functions as \(bird^c\). It operates within the inferential and generic relations the word maintains in \(\langle bird \rangle\) --- what is said about birds in the community, how the term combines with predicates, what counts as a counterexample. No particular bird is in view. The utterance is about how the word ``bird'' connects to ``fly'' within the linguistic pattern, with whatever generic force the form carries.

The two modes can be told apart by what speakers can and cannot do with them. A child who has acquired \(bird^o\) can point at a bird and say ``bird''; she can re-identify the same bird across short stretches of time. She cannot yet answer ``what is a bird?'' except by ostension. A child who has acquired \(bird^c\) can answer in words: ``a bird is an animal that flies.'' The difference is not in the word's reference but in the kind of work the word does --- coordinating with a particular versus moving through a network of other words. Both are linguistic capacities. Neither reduces to the other.

This last point bears emphasis because it is where the framework most often gets misread. The object-mode \(x^o\) is \emph{still language}. It is not a perception, not a piece of experience. Experience belongs to \(E\{X\}\) --- the activation, in a given encounter, of the perceptual-motor-affective constellation \(\{X\}\). The word \(bird^o\) does not \emph{consist of} the experience of the bird; it coordinates with it. That \(x^o\) is language rather than perception is not merely asserted: the two come apart in both directions. The experience can occur without the word --- \(E\{X\}\) is present in pre-linguistic infants and in animals with no word at all --- and the word, used in object-mode, can misfire on its object, coordinating with the wrong thing in the visual field or with something misperceived. A perceptual state is whatever it is; an object-mode use can be \emph{incorrect}, and correctable as such, which is the mark of a move in a linguistic practice rather than a perceptual or causal relation. (The objection that \(x^o\) is ``just perception'' is taken up in full in Chapter 8; what matters here is only that \(x^o\) is neither pure perception nor pure linguistic relation.) The two-side structure --- \(\langle x \rangle\) on the linguistic side, \(\{X\}\) on the experiential side, \(\mathcal{R}\) as the coordination between them --- keeps \(x^o\) from collapsing into either pure linguistic relation or pure perceptual content. A speaker using a word in \(x^o\)-mode is doing something linguistic, but doing it in a way that requires coordination with what is encountered.

The closest structural analogue in the existing literature is Recanati's mental files framework. The present account builds on similar insights but differs in a crucial respect: it locates the distinction at the level of publicly usable word-modes and social-linguistic practices rather than treating it primarily as a classification of private mental vehicles. Furthermore, it applies this distinction symmetrically to common nouns and names.

One refinement of the concept-mode is needed now, because it is relied on later. The generic \(x^c\) just described --- \(bird^c\) operating within \(\langle bird \rangle\) --- is not the only non-perceptual mode of a word. A proper name used in the bearer's absence is also non-perceptual, but it is not a generic concept-word of the same kind. Mark it \(N^c\): the name-mode use, a standing individual-tracker carried by a public name-pattern, whose function is to track one individual rather than to operate within the inferential and generic relations of a kind-term. A speaker who knows ``Aristotle'' only through books is not deploying a generic concept-word; she is deploying an inherited \(N^c\), supported by descriptive and testimonial material that mediates her access to the referent without constituting it. The distinction matters because it blocks a collapse the account must not allow: that every non-perceptual use of a name is a concept-mode use of the same type as \(bird^c\), so that a speaker in the bearer's absence operates ``solely at \(x^c\).'' She does not. She operates through an inherited name-pattern, \(N^c\), and need not repeat the original object-mode coordination that anchored the name. How \(N^c\) is anchored at baptism and carried across absence is the business of the next chapter; what matters here is only that name-mode is a distinct functional mode, not generic concept-mode wearing a proper name.

\subsection{\texorpdfstring{\(\mathcal{R}\) as Coordination}{\textbackslash mathcal\{R\} as Coordination}}\label{mathcalr-as-coordination}

The relation \(\mathcal{R}\) is the framework's coordination\footnote{This is a relation between two ways of using the same word, to describe that the concept word is grounded in a real experience. \(x^o\) and \(x^c\) are the same token or utterance ``x'', not two different tokens or utterances.} operator. It describes how the items of the two-side topology --- words in \(\langle x \rangle\) and patterns in \(\{X\}\) --- stand in relation to one another. The relation can be of several kinds, depending on what is being coordinated. A word in object-mode can coordinate with an experiential activation: \(\mathcal{R}(bird^o, E\{BIRD\})\) describes the binding by which ``bird'' used demonstratively coordinates with a particular bird in view. A word in object-mode can coordinate with a word in concept-mode: \(\mathcal{R}(bird^o, white^c)\) describes the predicative coordination by which ``that bird is white'' attaches the concept-mode predicate to the object-mode subject. A word in concept-mode can coordinate with another word in concept-mode: \(\mathcal{R}(bird^c, fly^c)\) describes the generic coordination by which ``birds fly'' connects the two concept-mode items within the linguistic pattern.\footnote{\(\mathcal{R}\) is specified in the framework document at greater technical length, including its developmental account in the learning of \(\{X\}\) and its role in the integrated \(\{X, x^o\}\) constellation produced by the coordination of experiential event and word; see PRU Meta-Language, version 6.}

In each case the operator names the same kind of phenomenon --- two items standing in coordination --- and the argument types determine the specific kind. The claim of the framework is that what looks like several distinct things in standard semantics --- reference, predication of a particular, generic predication, and baptism --- are applications of one relation. The claim earns its keep only if the unification does work that separate treatments do not, and it does: because baptism~\(\mathcal{R}(h^o,N^c)\)~and attribution~\(\mathcal{R}(\mathcal{D},N^c)\) are the \emph{same} operator differing only in argument type, the account can say precisely where reference-fixing and description part company (the former crosses the experience--language boundary, the latter stays within~\(⟨x⟩\)) ---a distinction the reply to descriptivism in the objections chapter turns on. A theory with four unrelated relations cannot locate the asymmetry in one place; this one can.

The word ``coordination'' carries the substantive load here. The standard alternative is Fregean function-application: a predicate is a function that takes an object as argument and yields a truth-value. Set-theoretic semantics produces a parallel architecture: a predicate is a set, predication is membership. Both pictures are asymmetric and arity-fixed: the predicate is the active item, the object the passive one, and the predication relation is the application of one to the other. The coordination picture differs at exactly this point. The two items in \(\mathcal{R}(x, y)\) stand in a structural relation that the practice sustains, not in a function-argument pairing the speaker computes. The asymmetry between ``subject'' and ``predicate'' in surface grammar is real but not constitutive: in \(\mathcal{R}(bird^c, fly^c)\) neither item is more ``argument'' than the other, and the same operator covers the case where one item is in object-mode and the other in concept-mode.

This is not a computational picture in an algorithmic or procedural sense.~R~is not an operation a speaker~\emph{performs}~on linguistic items, in the way a function gets applied to an argument or a set-membership relation gets checked. Rather,~\(\mathcal{R}\) functions as a scale-free formal mapping that tracks structural coordination across a cognitive continuum, spanning from sub-personal pattern-binding to socio-pragmatic embeddedness. To say that~\(\mathcal{R}(bird^o,white^c)\)~holds is to say that the speaker's use of ``white'' for this bird is sustained by an objective convergence across multiple tiers of this continuum. It tracks an ongoing engagement with both the bird (perceptually, at the level of object-mode invariants) and the predicate (within the linguistic space~\(⟨white⟩\), which governs application conditions, inferential roles, and contrast cases). This coordination is constituted developmentally---through the agent's transition from basic tracking mechanisms to public practices---and is sustained by continued participation in the linguistic community and continued contact with the world. The register is closer to Wittgenstein than to cognitive engineering, provided~\(\mathcal{R}\)~is understood as a scale-free formal mapping that tracks structural coordination across a cognitive continuum, rather than a computational procedure executed in real time. On this view, the high-level linguistic space~\(⟨x⟩\)~is formed when language games are learned inside a practice. For example,~\(⟨cat⟩\)~designates the totality of sentences using ``cat'' that make sense in a given language game, which~\(\mathcal{R}\) formally maps onto a corresponding practice with~\(\{CAT\}\), the lived cognitive attractor state.~\(\mathcal{R}\)~thus functions as a formal descriptor of how language is lived, structurally anchoring public linguistic practices to sub-personal connectionist dynamics without collapsing into anti-mentalism.

\subsection{\texorpdfstring{Local Rigidity at \(x^o\)}{Local Rigidity at x\^{}o}}\label{local-rigidity-at-xo}

Consider the utterance ``that bird could have been larger.'' A speaker pointing at a bird in front of her makes the counterfactual claim. The bird in question is held fixed across the modal variation: the speaker is not imagining a different bird that is larger; she is imagining \emph{that very bird} having been larger. The word ``bird'' in \(x^o\) mode coordinates with the same particular across the modal scope of ``could have been.'' Whatever properties of the bird are varied --- size, color, vigor, age --- the bird stays the same, by virtue of the demonstrative coordination already in place.

The modal test for rigid designation is whether a designator picks out the same individual across the worlds the modal context quantifies over.\footnote{Kripke, \emph{Naming and Necessity}, p.~48.} Within the modal scope of this utterance, the test passes: in every counterfactual scenario where ``that bird could have been larger'' is evaluated, ``that bird'' picks out the same bird that is in front of the speaker in the actual context. The variation occurs in what is predicated of the bird, not in which bird is being talked about. What the demonstrative coordination supplies directly is not rigidity but a grip on a particular --- the bird is held by the coordination \(\mathcal{R}(x^o, \cdot)\); rigidity is what that grip amounts to once the modal context is applied to it. To say the use is locally rigid is to say that the particular the coordination holds is what every counterfactual in the utterance is about.

Call this \emph{local rigidity}: rigidity that obtains within a perceptual or conversational episode by virtue of the demonstrative coordination \(\mathcal{R}(x^o, \cdot)\) being in place. Local rigidity is not yet Kripkean rigidity proper. Kripkean rigidity attaches to a name and holds across the entire range of counterfactual scenarios in which the bearer exists; local rigidity holds within the episode and lapses when the coordination does --- when the speaker stops attending to the bird, or when the conversational context shifts. The two differ in scope and social scaffolding, while exhibiting a closely related holding-fixed structure. What is varied is the durability of the coordination, not the cognitive operation that constitutes it.

A qualification is needed here, because a natural objection turns on it. Demonstratives are often said not to be rigid at all: their reference shifts with context, attention, and intention. This is correct of the demonstrative \emph{type}. ``That bird'' does not pick out the same bird across contexts; each use requires a fresh demonstration, and a different demonstration coordinates with a different particular. What is locally rigid is the demonstrative \emph{token}, once its demonstration is in place: given that this use coordinates with this bird, every counterfactual built on it concerns that bird and not another. The local rigidity claimed here is a claim about the token under a fixed coordination, not about the context-sensitive type --- which is why it suffices for the local \emph{de re} point without making demonstratives into portable rigid designators. Pure indexicals belong on the other side of this line. ``I,'' ``now,'' ``here'' are directly referential, but their referent is fixed automatically by a rule on the context parameters, with no demonstration and no encounter; they are not coordinations with a singled-out particular, and so they are not the kind of object-mode use at issue. The primitive here is demonstrative coordination specifically, not indexicality in general.

This last point bears emphasis, because the structural claim of the chapter depends on it. Rigidity is not a stiff property of the designator. It is a preserved link between the speaker's use of the term and the particular that anchors the use --- and the link, in the demonstrative case, is the coordination that grounds the use in an encountered particular. The preservation is constituted by the practice --- perceptual attention, in the demonstrative case; social-linguistic stabilization, in the case of names. To call a term rigid is to say that this link is being preserved across whatever variation the modal context introduces. The variation is in the properties imagined, not in the link itself.

It is worth setting this picture against the natural alternative. A different way of conceiving rigidity treats the designator as something with rigidity built in --- a kind of stiff attachment to its bearer that, once established, holds across modal variation by virtue of what the designator is. On that picture, what makes the modal evaluation come out right is the rigidity of the attachment itself, a feature the designator possesses. The picture is a natural one and is encouraged by some of Kripke's own formulations, but it misdescribes what is happening. Names do not have rigidity built in. Names rigidly designate because the linguistic community has a practice of using them rigidly --- of correcting failures of reference, of holding the term to its bearer across descriptive revision, of teaching the term so that successive speakers continue to use it in the same way.\footnote{Kripke's own metaphors are not in fact uniformly committed to a built-in-property reading. His description of the baptismal act and the causal chain (NN pp.~91--97) treats reference as something speakers do and sustain, which is closer to the preserved-link framing. The built-in-property reading is a tendency of the formal regimentation rather than of Kripke's positive picture; see Kripke's own characterization of his account as ``a better picture'' rather than a theory at NN p.~94.} The preserved-link framing captures what this practice accomplishes; the built-in-property framing obscures it.

Names, on the account developed here, are devices for sustaining the preserved link across the conditions that make demonstrative re-anchoring impossible: the bearer is not perceptually present, the speaker is not the baptizer, the chain of transmission is many speakers and centuries long. The cognitive coordination at work in \(\mathcal{R}(bird^o, \cdot)\) provides the model for the kind of coordination later sustained in absent-name use such as \(\mathcal{R}(\sigma, {N^c}_\text{aristotle})\) at the level of a surrogate-based reading of Aristotle two millennia after his death.\footnote{The notation \(\sigma\) for surrogate objects --- manuscripts, portraits, citations --- is from ``Pattern-Recognition Unity: A Framework Specification'' (Cojocariu, 2026).} What differs is the scaffolding that sustains the coordination: perceptual contact in the first case, testimonial networks and surrogate objects (manuscripts, portraits, citations) in the second. What is the same across the two is the holding-fixed structure --- a primitive coordination differently scaffolded; what differs is the practice that sustains it.

Names have been treated as indexicals/demonstratives before (Recanati, and others). Those are claims about names' semantic category; mine is a claim about the cognitive coordination any such semantics presupposes. So I don't compete with them --- I sit underneath them.

One clarification forestalls a natural misreading. The semantic literature already distinguishes kinds of rigidity --- obstinate from persistent rigidity (Salmon 1981),~\emph{de jure}~from~\emph{de facto}~(§ 3.1 above), the wide-scope ``rigidity'' a description can be given by scope placement. Local rigidity is not another entry in that taxonomy. Those are classifications of how a designator behaves across worlds once it is in hand; local rigidity is a claim about the coordination that any such designator presupposes when it is anchored to a particular at all. It sits below the taxonomy, not beside it --- the same relation to the rigidity literature that the account bears to Recanati's files at the level of singular thought.

The present account treats the demonstrative use of ``bird'' and the use of ``Aristotle'' as cases of closely related forms of coordination, differing in the practice that sustains the coordination rather than in the kind of cognitive operation involved.

\subsection{\texorpdfstring{\emph{De Re} and \emph{De Dicto}: Scope and the Holding-Fixed Operation}{De Re and De Dicto: Scope and the Holding-Fixed Operation}}\label{de-re-and-de-dicto-scope-and-the-holding-fixed-operation}

The \emph{de re} / \emph{de dicto} distinction is medieval in origin and gets one of its central modern formulations in Quine's puzzle about quantifying into modal contexts.\footnote{Quine, ``Reference and Modality,'' in \emph{From a Logical Point of View} (1953); Kaplan, ``Quantifying In,'' \emph{Synthese} 19 (1968): 178--214.} Consider the canonical example:

\begin{quote}
Necessarily, the number of planets is greater than seven.
\end{quote}

The sentence is ambiguous. On one reading --- \emph{de dicto}, or narrow scope --- ``the number of planets'' is evaluated at each world: in every possible world, the number of planets in that world must be greater than seven. The reading is false. On the other reading --- \emph{de re}, or wide scope --- the description first picks out the number that is actually the number of planets, and the modal claim is then made about that number. Of that number, namely eight, it is necessarily true that it is greater than seven. The reading is true.

Two standard treatments handle this. Russell's scope-of-the-quantifier solution treats the ambiguity as one of logical form: the description can take scope inside or outside the modal operator.\footnote{Russell, ``On Denoting,'' \emph{Mind} 14 (1905): 479--93, on scope generally; the modal application is developed in Smullyan, ``Modality and Description,'' \emph{Journal of Symbolic Logic} 13 (1948): 31--37.} Kripke's rigidity solution treats the relevant contrast as a feature of the difference between names and descriptions: names are rigid and behave like wide-scope descriptions in modal contexts; descriptions can be evaluated either way.\footnote{Kripke, \emph{Naming and Necessity}, pp.~48--60 and the scope footnote at p.~12.} The present account does not replace these treatments. It adds a further question: what cognitive-linguistic operation allows a speaker to hold an object fixed for \emph{de re} modal predication in the first place?

This is where the \(x^o/x^c\) distinction enters. It should not be read as a rival to scope theory. A reader may hold that scope theory already handles \emph{de re} predication completely --- that once the description's scope relative to the modal operator is fixed, there is nothing further to explain, and the question of what lets a speaker hold an object fixed is not a question about reference but about psychology. The reply is that scope theory tells us how the two readings differ in logical form but is silent on what a speaker does to make the wide-scope, object-fixed reading available at all; that doing is the present topic, and it is left untouched by the scope account rather than competing with it. In the ``number of planets'' case, the fixing is not perceptual R-grounding. Numbers are not encountered particulars in the way birds, tables, or persons are. The object is secured by a formal or descriptive route, so the case is A-grounded rather than R-grounded. It is structurally analogous to holding something fixed, but it is not the paradigm case of \(x^o\)-coordination.

This concession should be registered explicitly. Holding a target fixed for de re predication has two anchor types: R-grounded, where the target is secured in encounter, and A-grounded, where it is secured within \(⟨x⟩\)-space. The dissertation's thesis concerns the first: the claim is that the rigidity of~\emph{names of encountered particulars}~derives from local rigidity, not that every instance of holding-fixed does. Whether the two anchor types are exercises of a single capacity is taken up in the concluding chapter.

The clearer case is demonstrative:

\begin{quote}
That bird could have failed to fly.
\end{quote}

Here the speaker is not saying that, in some possible world, whatever satisfies the description ``bird in front of me'' fails to fly. Nor is she making a general modal claim about birdhood. She is holding fixed this very bird and varying what is predicated of it. In the vocabulary of this dissertation, the use of ``that bird'' is anchored through \(R(bird^o, E(\{BIRD_i\}))\), where \(E(\{BIRD_i\})\) is the active encounter with that particular bird. This is the local form of \emph{de re} predication: the particular is fixed first; modal predication proceeds afterward.

Proper names can be understood as a socially stabilized extension of this structure under conditions of absence. Once a name has been successfully introduced and transmitted, its ordinary modal use resembles the \emph{de re} side of the contrast: ``Aristotle could have failed to teach Alexander'' holds Aristotle fixed and varies what is predicated of him. But the current speaker does not perform fresh R-grounding with Aristotle. The name functions \emph{de re} because a linguistic practice preserves the reference fixed by earlier acts of naming and transmission.

This also clarifies Quine's worry about essentialism.\footnote{Quine, ``Reference and Modality'' (1953); ``Three Grades of Modal Involvement,'' \emph{Proceedings of the XIth International Congress of Philosophy} 14 (1953): 65--81.} The account does not dissolve the metaphysical issue. Holding an object fixed explains how \emph{de re} modal predication is possible as a cognitive-linguistic practice; it does not decide which properties, if any, the object has essentially. It separates two questions: how a speaker can make a modal claim about a particular as that particular, and which modal claims about that particular are true.

The contribution of the \(x^o/x^c\) distinction is therefore modest but important. Scope theory tells us how the readings differ in logical form. The object-mode/concept-mode distinction explains how speakers are able to hold a particular fixed in the first place, and why demonstrative and name-based reference behave differently from purely descriptive or conceptual deployment.

\subsection{\texorpdfstring{The Necessary \emph{A Posteriori} at the Level Distinction}{The Necessary A Posteriori at the Level Distinction}}\label{the-necessary-a-posteriori-at-the-level-distinction}

Kripke's claim that some identities are necessary but \emph{a posteriori} broke a connection that had been taken for granted in much of the analytic tradition: that necessity goes with analyticity and \emph{a priori} knowability, while \emph{a posteriori} truths are contingent. The classical examples are identities involving rigid designators --- ``Hesperus is Phosphorus,'' ``water is \(H_2O\),''---as well as necessary a posteriori predications such as ``this very table is made of wood. Each is necessary, on Kripke's account, because both flanking terms rigidly designate the same particular or natural kind. Each is \emph{a posteriori} because the identity was established empirically.\footnote{Kripke, \emph{Naming and Necessity}, pp.~100--105 on Hesperus/Phosphorus; pp.~116--18 on water/\(H_2O\); the general thesis is at pp.~35--39. The pre-Kripkean connection between necessity, the \emph{a priori}, and analyticity was defended in different forms by Carnap, Ayer, and others, and is the standard target of Kripke's first lecture.}

The phenomenon has been read by some as paradoxical. If ``Hesperus is Phosphorus'' is necessary, why was the identity not known \emph{a priori}, like ``2 + 2 = 4''? If it is \emph{a posteriori}, how can it be necessary rather than contingent? The puzzlement is real on a picture where necessity is identified with analyticity and epistemic accessibility tracks metaphysical modality. On the picture developed in this chapter, the puzzlement has a clean source: it conflates two levels of language use that should be kept apart.

\(x^o\) supplies the fixed object whose identity grounds the necessity; \(x^c\) explains why recognition of that identity is empirical. When Hesperus is observed in the evening and Phosphorus in the morning, the same celestial body is the object of two demonstrative coordinations --- once at dusk, once at dawn. At the level of \(x^o\), there is only one Venus: \(E(\{VENUS\})\) activated in the evening is, as it happens, \(E(\{VENUS\})\) activated in the morning. The identity of the planet is constitutive at the demonstrative-coordination level: the same particular is being tracked in both cases. If the evening and morning observations do in fact concern the same celestial body, then the relevant identity lies at the level of the object being tracked, not at the level of the names or descriptions through which speakers conceptualize it. The empirical discovery is precisely the realization that two independently established modes of access actually converge on one and the same object.

The \emph{a posteriority} lives at \(x^c\). The concept-mode designators \(\langle Hesperus \rangle\) and \(\langle Phosphorus \rangle\) are distinct linguistic patterns. Each was established by an independent baptismal stabilization: one anchored in the evening-sky coordination, the other in the morning-sky coordination. As concept-mode terms they have different usage histories, different inferential connections, different roles in their respective linguistic patterns. Establishing that they coordinate to the same particular requires empirical work --- astronomical observation, theoretical inference, and acceptance of the identification by the community. The inference \(\mathcal{R}(Hesperus^c, Phosphorus^c)\), that the two concept-mode patterns share a referent, is \emph{a posteriori} in the standard sense, where this coordination is read as co-anchoring, not as the generic coordination of §4.2.

The phenomenon is not a paradox on this reading. It is a structural fact about how language operates at two levels. The identity is necessary at the level where the same particular is being coordinated with; the discovery of the identity is \emph{a posteriori} at the level where the two linguistic patterns are reconciled. The statement is one statement, but its modal and epistemic profiles are explained by different aspects of our practice: the fixed referent explains necessity; the historical/conceptual route by which speakers discover co-reference explains a posteriority. The traditional connection between necessity and \emph{a priori} knowability assumed that there is only one level at which the statement operates. The \(x^o/x^c\) distinction blocks the assumption.

The same applies to ``water is \(H_2O\).'' At the level of \(x^o\), the natural kind whose instances were demonstratively tracked by speakers using ``water'' is the same natural kind whose instances are chemically composed of \(H_2O\) molecules. There is one kind, and the demonstrative tracking by ordinary speakers and the chemical analysis by scientists coordinate to it. The identity is necessary at this level. At the level of \(x^c\), the two terms --- the ordinary-language \(\langle water \rangle\) and the chemical-formula \(\langle H_2O \rangle\) --- are different linguistic patterns with different inferential roles. Establishing that they coordinate to the same kind required scientific investigation. The identity is \emph{a posteriori} at this level.\footnote{The water-\(H_2O\) case raises additional questions about natural kinds, essence, and the role of scientific theorizing in fixing reference. These are discussed in Putnam, ``The Meaning of `Meaning','' in \emph{Mind, Language and Reality} (1975), and have generated a substantial literature. The present section makes no commitment on the metaphysics of natural kinds; the level distinction handles the modal-epistemic shape of the case without taking a position on whether being \(H_2O\) is the essence of water in any deeper sense.}

The necessary \emph{a posteriori} has been handled in the literature through theories of trans-world identity, essentialism about substances and natural kinds, and two-dimensional semantics.\footnote{Salmon, \emph{Reference and Essence} (1981), and the survey in Soames, \emph{Beyond Rigidity} (2002), chs.~7--9.} The present treatment does not contest any of these accounts at the semantic level. It locates the apparent paradox at the cognitive-linguistic level: necessity at \(x^o\), \emph{a posteriority} at \(x^c\). The two levels were conflated in the tradition Kripke broke from, and they remain easy to conflate even after Kripke. The level distinction marks them apart and gives a structural account of why the necessary-\emph{a posteriori} cases have the shape they have.

\section{The PRU Account of Rigid Designation}\label{the-pru-account-of-rigid-designation}

The present chapter offers the constructive proposal. It does not contest Kripke's semantic results; it explains how speakers come to inhabit the models in which those results hold. The argument runs from \(x^o\)-mode tracking, through baptismal stabilization, through testimonial transmission, to the description-attribution practices by which most of us actually engage with names whose bearers we have never encountered, and ends with the metasemantic positioning of PRU relative to standard Kripkean model theory.

\subsection{The PRU Framework}\label{the-pru-framework}

Only what is needed for the argument will be introduced here. The full framework is developed in a distinct preprint (Cojocariu, 2026). Three notions matter.

A \textbf{pattern-constellation} \(\{X\}\) is a unified attractor state in a cognitive system, sculpted by repeated encounters and integrating sensory, motor, affective, and interoceptive dimensions as a single learned configuration. We never encounter \(\{X\}\) as such; we encounter particular instances. A pattern-recognition event \(E(\{X\})\) is the system's falling into the \(\{X\}\)-basin in response to a partial cue. Recognition and reaction are aspects of one event, not a detection step followed by a separate response --- this is the Pattern-Recognition Unity thesis from which the framework takes its name.

A \textbf{linguistic pattern-constellation} \(\langle x \rangle\) is the analogous attractor in linguistic space --- the standing pattern of uses, collocations, inferential connections, and discursive habits associated with a word. \(\langle x \rangle\) is to language what \(\{X\}\) is to experience.

Words can be employed in two different modes. As object-word, noted \(x^o\), is learned as a motor pattern inside a constellation pattern \(\{X\}\). At this stage the infant says ``cat'' when experiencing \(\{CAT\}\) but cannot explain what a cat is. Later on, immersed in the language game of adults, the kid learns uses of ``cat'' in sentences when there is no cat present. This use is called concept-word, noted \(x^c\). When asked what a cat is, the kid can say ``A cat is a furry animal'', using \(cat^c\). We only use, as adults, \(cat^o\), in the presence of the designated object, like saying ``cat!'' or ``this cat'', often with a pointing gesture indicating the real cat. This is also the word ``appearing'' mentally in an involuntary fashion when an object is perceived. Any use of ``cat'' in a fully formed sentence that does not refer to a real, actual, perceived cat, is in fact \(cat^c\) use.

The \textbf{coordination operator} \(\mathcal{R}\) couples these two kinds of constellation. At the surface level we write \(\mathcal{R}(x^o, x^c)\) --- relating the word in its object-mode use to its concept-mode entry\footnote{This is a relation between two ways of using the same word, to describe that the concept word is grounded in a real experience. \(x^o\) and \(x^c\) are the same token or utterance ``x'', not two different tokens or utterances.}. At the deeper architectural level the same operation is \(\mathcal{R}(\{X\}, \langle x \rangle)\) --- coupling the experiential basin to the linguistic pattern. The two notations refer to the same relation under different resolutions. As noted when the operator was introduced, ``coordination'' is retained as the working term even though, in its world-crossing uses, what it describes may be at bottom a grounding relation; throughout, the coordination at issue is the ``in the moment'' relation between a word and what it is used about, not the ``in the time'' alignment of speakers with one another.

That is the apparatus. The remainder of this chapter is what it does.

\subsection{\texorpdfstring{Local Rigidity: \(x^o\) as the Language Primitive}{Local Rigidity: x\^{}o as the Language Primitive}}\label{local-rigidity-xo-as-the-language-primitive}

Russell's acquaintance was meant to identify a special form of reference: a contact with an object that secures identification independent of any descriptive material. He restricted it to immediate sense-data because only there did he find the Cartesian certainty he took such reference to require. Once that restriction is dropped --- once we allow that reference can be stable without being incorrigible --- Russell's mechanism generalizes. What he saw in the sense-datum case is what happens whenever a speaker tracks a particular through encounter: a word in \(x^o\)-mode coordinates with an experiential instance, and that coordination grounds the use in the particular. The coordination is what is locally rigid --- rigid in the sense that, for as long as the tracking holds, every counterfactual built on the use is about that particular. ``That bird could have been larger'' requires me to hold the bird fixed across the modal evaluation; the holding is accomplished by my present perceptual contact, not by any description I could supply.

\(x^o\) is the PRU implementation of acquaintance with two specific modifications. PRU's grounding is \textbf{fallible}: \(E(\{X\})\) can misfire, the basin \(\{X\}\) is reshaped by later encounters, and I can be wrong about properties of what I am tracking. PRU's grounding is also \textbf{dynamic}: the attractor state is sculpted continuously, not given once and for all. These features make \(x^o\) less like Russell's punctual sense-datum and more like the standing capacity of a speaker to keep a particular fixed through engagement with it. Kripke's own account already permits the speaker to be wrong about the bearer of a name; the rigidity does not depend on infallibility. PRU pushes this point further into the cognitive substrate: tracking does not require certainty, only stability.

The cognitive primitive is therefore the \(x^o\)-coordination---the word in object-mode locked onto an experienced instance. At the primitive level, this coordination exhibits the clearest case of non-descriptive holding-fixed; it provides the basic model from which the later account proceeds. It is not a property added to reference; it is what reference is at its primitive level.

A word on the order of explanation the account assumes, promised in its broad shape earlier. The demonstrative coordination is what is given first; rigidity is not a further ingredient alongside it but the shape that coordination takes once a modal context is brought to bear. When a speaker holds a particular through \(x^o\)-coordination, there is nothing in the use for a modal operator to vary except what is predicated of the particular --- and that is just what it is for the use to be locally rigid. On this account, then, the grounding of a word in a particular comes first, and rigidity is its modal manifestation; the chapter develops the consequences of taking the coordination, rather than rigidity, as the primitive. This is the order of dependence the rest of the account runs on. It is stated here as the account's framing, not argued as a thesis: what grounding is at bottom, and whether rigidity reduces to it or merely manifests it, are the deeper questions flagged earlier and left to separate work. What the account needs, and uses, is only the directional claim --- coordination first, rigidity as its modal face --- which the construction that follows puts to work.

\subsection{The Problem of Absence}\label{the-problem-of-absence}

The primitive form of reference has a limitation built into it. When the bearer leaves the perceptual field, the \(x^o\)-coordination dissolves. ``That bird'' stops functioning when the bird flies away. What had been a rigid grip becomes just a token.

The natural descriptivist response is to substitute description for tracking: in the absence of the bird, refer to it as ``the bird I saw earlier in the willow.'' But the chapter on rigid designation has already shown that descriptions cannot do the work of names. They pick out world-by-world satisfiers; they fail in modal contexts; they are constitutively \(x^c\) and have no \(x^o\) anchoring of their own to draw on. A speaker who has only descriptions to work with has, in the relevant sense, lost the particular.

This is the problem names solve. A name is a device for extending, under social conditions, the holding-fixed function first visible in \(x^o\)-coordination into the space of absence. The grip is real where the object is present to the speaker; what is needed is a way to carry it forward when the object is not present to the speaker. Names are best understood as portable devices for maintaining singular reference: linguistic artifacts that keep a word anchored to one individual when perceptual contact is no longer available. The word ``portable'' should be heard with care. What a name carries forward is the singular anchoring --- the fixing of one individual --- not a ready-made modal property handed over from perception. In what sense the resulting name counts as rigid, and how the anchoring survives the loss of contact, is the work of the sections that follow.

\subsection{Baptism as Stabilization}\label{baptism-as-stabilization}

A baptismal act exploits an existing \(x^o\)-coordination to inaugurate a new linguistic pattern. At the moment of baptism, some speaker is in perceptual contact with the individual to be named. \(\mathcal{R}(human^o, E(\{HUMAN\}))\) is in place, I experience a particular human, this one I can call ``this human''; the object-mode word ``human,'' or ``child,'' or ``this one,'' is rigidly coordinated with the particular instance. The baptism does two things simultaneously: it introduces a new concept-word \(N^c\) (the name), and it forms a coordination between the already-rigid \(h^o\) and this newly introduced \(N^c\):

\[\mathcal{R}(h^o, N^c)\] for example: \[\mathcal{R}(h^o, {N^c}_{aristotle})\] This effectively states, `Let this human be named Aristotle.' But because proper names are introduced as standing linguistic devices rather than as passing demonstratives, what this accomplishes is practical rather than metaphysical. A new linguistic pattern \(\langle aristotle \rangle\) enters the community's vocabulary. Its standing function --- fixed by the baptismal context and maintained by subsequent use --- is to track the same individual the demonstrative grip was holding. What baptism does is \textbf{stabilize} the grounded fixation: the live \(x^o\)-coordination held one particular for the duration of contact, and baptism gives that same singular fixation a durable, publicly maintained vehicle in \(N^c\). The change is one of \emph{status} --- from a fixation confined to the duration of perceptual contact to one available for use beyond it --- and not the transfer of a ready-made property from perception into the name.

This is why it is not baptism itself that originates the singular fixation. What was already in place, in \(\mathcal{R}(h^o, E\{HUMAN\})\), was a grip on one particular --- grounded in the encounter, holding that individual and no other. What baptism changes is the locus of that grip: from a fleeting demonstrative coordination to a durable, socially maintained linguistic pattern. The name can then be treated, in Kripke's sense, as a rigid designator --- and the reason it can is structural. A description has descriptive content that interacts with modal operators, generating the scope ambiguities discussed in the chapter on the functional-mode distinction; \(N^c\) has no such content to interact with. Anchored to its individual and carrying no condition that could pick out different satisfiers in different worlds, it has nothing for a modal operator to vary. Its rigidity is the modal face of a fixation that was settled at the encounter and carries no descriptive structure --- not a separate property conferred by the naming.

Kripke's ``initial baptism'' is the same event described from the outside. What his account leaves implicit --- that some speaker had to be tracking the individual at the moment the name was introduced --- is what PRU makes explicit. The picture is not new; the cognitive precondition is.

The account as stated covers baptisms in which the baptizer is in perceptual contact with the bearer --- ostensive baptism, where the founding coordination is \(\mathcal{R}(h^o, N^c)\), crossing from an encountered particular to the name. Baptisms that fix reference through a description of an unencountered cause --- ``Neptune,'' ``Jack the Ripper'' --- have a different structure: there is no founding \(x^o\)-coordination with the bearer, and the name is introduced by a coordination of the type \(\mathcal{R}(\mathcal{D}, N^c)\), running from a description to the name within \(\langle x \rangle\)-space. Whether such a baptism anchors its name in the way ostensive baptism does, or anchors it only through the naming norm, is a question the ostensive case does not settle. Naming the contrast is all that is needed here; the descriptive and deferred cases are taken up among the open questions of the concluding chapter.

\subsection{Causal Chains as Testimonial Networks}\label{causal-chains-as-testimonial-networks}

The question is how rigid tracking propagates through a community when only the first speaker(s) had \(x^o\)-access. What is the mechanism of the causal chains that Kripke mentions?

It helps to say first what ``tracking'' means once the bearer is absent, because it is no longer the perceptual tracking of §5.2. When a contemporary speaker tracks Aristotle, she is not in perceptual contact with him and performs no \(x^o\)-coordination with him; there is no live grip to maintain. What she maintains is a \emph{use}: her continued, norm-governed deployment of an inherited name-pattern \(\langle aristotle \rangle\) that the community keeps anchored to one individual. Tracking across absence is therefore not a cognitive relation between the speaker and the bearer but the speaker's sustained participation in a public practice that holds the name to its origin --- a practice constituted by training, correction, and the testimonial material through which the name is transmitted. The grip on the particular was established once, at the encounter that founded the name; what later speakers do is keep using the name under the norm that preserves that anchoring, not re-establish contact with the individual. The channels below are what make this sustained use possible.

Three transmission channels operate. \textbf{Linguistic patterns} \(\langle N \rangle\) carry the bulk of the work: a speaker who has heard ``Aristotle'' used in the standard ways forms her own \(\langle aristotle \rangle\) --- a stable pattern of uses, collocations, and inferential connections that mirrors the community's. \textbf{Testimonial patterns} \(\tau(N)\) are the statements, texts, narratives, and corrections that flow through the network: an academic conference, a textbook, a teacher's correction of a student's misuse. These shape and stabilize \(\langle N \rangle\) for the speaker who internalizes them. \textbf{Surrogate constellations} \(\sigma(N)\) are the residual experiential anchors: manuscripts, busts, recordings, photographs --- objects that causally trace to the original bearer and provide a partial \(E(\{X\})\) for speakers who encounter them. Surrogates are not necessary for reference. A blind historian uses ``Aristotle'' correctly without ever encountering an image. But for most speakers, surrogates supply the residual phenomenological connection that pure testimony cannot.

What is transmitted through these channels is~\(N^c\)~--- not the bearer, and not any particular description. The name-concept retains its standing function because the community treats it as a tracking device rather than a classifier. This is the practice in which rule-following becomes intelligible (Wittgenstein, \emph{Philosophical Investigations}, §§198--202): what the community's collective treatment of~\(\langle N \rangle\)~sustains is the \emph{persistence} of the name's anchoring across absence --- its continuing to hold to the individual fixed at the founding encounter --- kept in force by training, correction, and institutional embedding. The practice does not constitute the rigidity; the rigidity is the modal face of an anchoring established when the name was founded. What the practice constitutes is that the anchoring survives the loss of contact, carried across speakers and generations who never had it.

This is what Kripke's causal-historical chain comes to in cognitive terms. The chain is not a metaphysical wire between baptism and current use; it is a pattern of practices that keeps \(N^c\) doing its tracking work across generations.

It is worth being explicit about where this leaves the coordination between speakers, since that is the phenomenon a different tradition takes as central. The transmission just described is coordination in the social sense --- speakers aligning on a referent, deferring to one another, converging through correction and shared use. The present account does not deny this or compete with it; the channels above are a description of it. The claim is one of order, not of exclusion. Social coordination operates on a name whose anchoring is already in place: baptism grounds the name in an individual, and the community then coordinates on the name so grounded. The relational facts a relationist account locates --- co-reference sustained across a practice --- are real and have their home here, in the transmission layer. What the account does not grant is that these facts are where reference is \emph{fixed}. They are how a fixed anchoring is kept in circulation, not what establishes it. The between-speakers question and the word-to-particular question are both genuine; the order of explanation runs from the second to the first, which is why the social phenomena can be accommodated without taking inter-speaker coordination as the ground of reference. Whether that ordering is finally correct against a developed relationism is a larger question, and it is left to separate work; what matters here is that the account has a place for the social facts without resting on them.

A reader may press the opposite order. If a name's reference is kept in place by a community's practice of transmission and correction, why not say that the practice is what reference \emph{consists in} --- that the talk of a founding coordination is an idle wheel, and that transmission is the whole story? The account's answer is that transmission preserves an anchoring it does not create: a practice of keeping a name attached to ``the same individual'' presupposes that there is an individual the name was attached to, which is what the founding coordination supplies and what a chain with no origin would lack. But the disagreement is real, and a reader who holds that the practice needs no founding grip --- that origins drop out and only the ongoing use matters --- rejects the account at this point rather than later.

\subsection{Attribution after Transmission}\label{attribution-after-transmission}

A downstream speaker who acquires \(N^c\) through the causal chain---but who lacks direct perceptual access to the bearer---is in an asymmetric position. The name in her vocabulary inherits its rigidity from a baptismal coordination she never performed. Her cognitive engagement with the bearer must rely on a different mechanism. What is she actually doing when she uses the name?

The classical descriptivist answer is that she uses the name as shorthand for the descriptions she associates with it. As Kripke demonstrated, this fails because it leaves reference hostage to the accuracy of those descriptions. The present framework offers a structurally different answer: descriptions do not constitute reference; rather, they are attributed to a name-concept that already carries a fixed reference.

Let \(\mathcal{D}\) denote a description, formally an \(\iota\)-structured compound within \(\langle x \rangle\)-space. For Aristotle, a contemporary scholar might form several such descriptions:

\[\mathcal{D}_{1} := \iota x\,(x\text{ taught Alexander the Great})\] \[\mathcal{D}_{2} := \iota x\,(x\text{ wrote the }\textit{Nicomachean Ethics})\] \[\mathcal{D}_{3} := \iota x\,(x\text{ studied under Plato})\]

Each of these is a \(\langle x \rangle\)-side object: a linguistic pattern with a uniqueness presupposition, evaluable at any world to yield a satisfier (or to fail). The scholar's cognitive engagement with Aristotle consists in coordinating these descriptions with the inherited name-concept:

\[\mathcal{R}(\mathcal{D}_{1}, N^c_{Aristotle}) \wedge \mathcal{R}(\mathcal{D}_{2}, N^c_{Aristotle}) \wedge \mathcal{R}(\mathcal{D}_{3}, N^c_{Aristotle}) \wedge \ldots\]

Each conjunct is an \textbf{attribution}: a coordination running from a description to a name-concept. Several features of this structure matter.

The operator is the same \(\mathcal{R}\) used throughout. There is no separate ``attribution operator.'' What distinguishes attribution from baptismal grounding is the type of its arguments. Baptismal grounding is \(\mathcal{R}(h^o, N^c)\) --- an \(x^o\)-side argument coordinated with a \(\langle x \rangle\)-side argument, crossing the experience-language boundary. Attribution is \(\mathcal{R}(\mathcal{D}_i, N^c)\) --- both arguments live in \(\langle x \rangle\)-space, and the coordination stays within language. The operator's form is constant; what changes is the work it does, depending on what it is applied to.

The conjunction is the right form. Each attribution is independently established and independently revisable. Discovering that bifocals were not in fact invented by Franklin requires me to drop \(\mathcal{R}(\mathcal{D}_{\text{bifocals}}, N^c_{BF})\). The other attributions --- that Franklin signed the Declaration, that he was Postmaster General --- remain intact. My reference to Franklin is not disturbed by the revision because reference was never carried by any attribution. Reference is carried by \(N^c\), anchored at baptism by \(\mathcal{R}(h^o, N^c_{BF})\), propagated through the chain. The attributions are how I cognitively inhabit the name, not how the name secures its bearer.

For convenience we may write the shorthand \(\mathcal{R}(\mathcal{D}, N^c)\) when the internal structure does not need to be explicit. This abbreviation must be read with care: it stands for a conjunction whose individual conjuncts can be dropped, replaced, or revised without invalidating the others. The shorthand \(\mathcal{D}\) is a structured collection, not a monolithic description.

Two consequences follow.

The asymmetry between rigid names and non-rigid descriptions becomes formally visible. \(N^c_{Aristotle}\) is rigid in virtue of the baptismal \(\mathcal{R}(h^o, N^c_{Aristotle})\) it inherits. Any \(\mathcal{D}_i\) remains non-rigid: it picks out world-by-world satisfiers, and in worlds where Aristotle did not study under Plato, \(\mathcal{D}_3\) either fails or picks out someone else. The attributive \(\mathcal{R}(\mathcal{D}_3, N^c_{Aristotle})\) is itself a contingent claim about the actual world. Each attribution can be false. None of these contingencies touches the rigidity of \(N^c_{Aristotle}\).

The Hesperus/Phosphorus case acquires a transparent structure. Before the astronomical discovery, two name-concepts had been independently established --- two baptisms of one celestial body, one anchored in the evening-sky coordination, one in the morning-sky coordination --- and independently propagated, each with its own attributions: \[\mathcal{R}(\mathcal{D}_{\text{morning star}}, N^{c}_{\text{Phosphorus}})\] \[\mathcal{R}(\mathcal{D}_{\text{evening star}}, N^{c}_{\text{Hesperus}})\]

The discovery is the recognition that the two chains are anchored in the same individual:

\[\mathcal{R}(\textit{star}^{\,o}, N^{c}_{\text{Hesperus}}) \;\wedge\;
\mathcal{R}(\textit{star}^{\,o}, N^{c}_{\text{Phosphorus}})\]

Note what the discovery is \emph{not}: it is not the identity of the name-concepts. \(\langle\text{Hesperus}\rangle\) and \(\langle\text{Phosphorus}\rangle\) remain distinct linguistic patterns, with distinct usage histories and distinct attribution-sets --- which is why the identity statement is informative, and why its a posteriority has a subject matter (the necessary a posteriori, treated earlier). What is one is the anchor; what were two, and remain two, are the concept-words. Co-anchoring, not concept-identity, is the formal content of ``Hesperus is Phosphorus.'' The community's subsequent preference for a third name, ``Venus,'' is a fact about the practice --- chains carrying co-anchored names tend to consolidate --- not a logical consequence of the discovery.

The necessary a posteriori falls out structurally: necessity comes from the~\(\mathcal{R}(x^o,N^c)\)~stratum; a-posteriority from the~\(\mathcal{R}(\mathcal{D},N^c)\)~stratum; both strata are present in any contentful use of the name.

This is also what allows the cognitive engagement of a contemporary scholar with a long-dead figure to be substantive without being descriptivist. The scholar engages with Aristotle through her attributions. The attributions can be revised, refined, defended, or dropped. Through all such revisions, she continues to refer to the same individual --- because \(N^c_{\text{Aristotle}}\) continues to do its rigid work, independent of what she attributes to it.

\subsection{PRU as Metasemantics}\label{pru-as-metasemantics}

The standard semantics for modal logic assigns each name the same individual at every world; that world-constant assignment is what a name's rigidity comes to within the formal theory, and it is written in by stipulation. The semantics does not say how a real speaker comes to use a name in conformity with it.

PRU does not propose a different model. It does not contest the formal semantics of rigid designation. What it offers is an account of how speakers come to inhabit Kripkean interpretations: how the rigidity that the model stipulates as a formal constraint is achieved in cognitive practice. The developmental trajectory has four stages:

The \(x^o\)-coordination is the cognitive primitive, present already in pre-linguistic tracking and in the demonstrative uses of common nouns. The baptismal stabilization \(\mathcal{R}(h^o, N^c)\) is the act by which a name-concept takes over the singular anchoring of its founding demonstrative grip --- the fixing of one individual, given a durable public vehicle. The testimonial transmission carries \(N^c\) across speakers and generations without requiring further \(x^o\)-access. The attributive coordination \(\mathcal{R}(\mathcal{D}, N^c)\) is the cognitive mode in which most speakers, most of the time, engage with names whose bearers are absent or have always been absent.

The model stipulates that~the name has its referent fixed~and world-constant; PRU~explains how a speaker is~positioned to use the name such that~the stipulation correctly~describes her semantic practice.

This is the sense in which PRU is metasemantics, not rival semantics. PRU supplies what Kripke's apparatus presupposes --- the cognitive precondition Russell located and retreated from --- without contesting the formal semantics.

The argument is now in place. The objections that this account is bound to attract are taken up in the next chapter.

\section{Objections and Replies}\label{objections-and-replies}

The account is now in place: rigid designation is the social stabilization of local rigidity, baptism stabilizes an \(x^o\)-coordination into a portable \(N^c\), and contemporary speakers engage absent bearers through revisable attributions \(\mathcal{R}(\mathcal{D}, N^c)\) that ride on a reference the attributions do not constitute. An account of this shape attracts objections of three kinds. The first says it has not added anything: that the A-/R-grounding distinction relabels a contrast already drawn. The second says it has added the wrong thing: that for all the anti-descriptivist language, the treatment of absent reference is descriptivism with extra notation. The third says it has helped itself to what it needed to prove: that the persistence of reference across absence rests on an unargued categorial feature of names. A fourth, more local, presses on the claim that \(x^o\) is a mode of language rather than a relabelled perceptual relation --- the claim on which the whole scaling-up of demonstrative rigidity depends. I take them in that order. None is a strawman; the first three are the objections I would raise against this account if someone else had written it, and the fourth is the one a careful reader reaches by the end of the constructive chapter.

\subsection{The Restatement Objection}\label{the-restatement-objection}

The objection. The distinction between A-grounding and R-grounding, the critic says, is the old distinction between reference-by-description and reference-by-acquaintance wearing new clothes. Russell drew it. Evans redrew it as the contrast between descriptive and demonstrative thought. Recanati's typology of files is built on it. Renaming the descriptive side ``A-grounding'' and the demonstrative side ``R-grounding,'' and the demonstrative-rigidity phenomenon ``local rigidity,'' does not advance the debate; it supplies vocabulary. The thesis, on this reading, is a notational variant of a position already on the table, and its appearance of novelty comes entirely from the new symbols.

Reply. The objection would land if the contribution were the distinction. It is not. The acquaintance/description contrast is indeed old, and the dissertation says so: Russell drew it and then misplaced its boundary. What the account adds is not a new name for the contrast but a claim about its direction of explanation --- and that claim is substantive, falsifiable, and absent from the views the objection cites.

The claim is that rigidity flows upward from demonstrative coordination to names, rather than belonging to names in the first place and being approximated by demonstratives. On the acquaintance tradition, the demonstrative case and the name case are two species of singular reference, sitting side by side; acquaintance is one route to a singular thought, description another. On the present account they are not coordinate species. Local rigidity is the primitive, and Kripkean rigidity is what local rigidity becomes when a social practice carries it past the encounter. This is a claim of derivation, not of classification, and it has consequences the classificatory view does not have.

Two of those consequences are testable against the rival's own commitments. First, on the derivation claim, common nouns used demonstratively (``that bird'') must already exhibit rigidity, because the rigidity of names is built from exactly that material. The acquaintance tradition, focused on singular thought about individuals, has no reason to expect rigidity at the level of a kind-term used to pick out a passing instance; the present account requires it. Second, on the derivation claim, the failure mode of a name should be a failure of the practice that sustains the stabilized coordination --- a broken chain, a corrupted transmission --- and not a failure of the speaker's descriptive knowledge. That is precisely the Feynman and Gödel data Kripke used against descriptivism, now predicted rather than merely accommodated. A relabelling predicts nothing. This account predicts both, and would be refuted if demonstrative kind-reference turned out to admit scope ambiguity, or if reference tracked descriptive competence rather than chain-membership.

So the distinction is old and the dissertation neither claims nor needs otherwise. What is new is the order of priority among its terms, and the order of priority does work the side-by-side picture cannot do.

There is a sharper form of the objection, aimed not at the acquaintance tradition in general but at Recanati's mental files in particular, since that account comes closest to the present one. Recanati, too, makes the reference of a singular vehicle turn on a non-descriptive relation rather than on satisfied descriptions, and houses the descriptive material in a file that can be wrong about its object without ceasing to be about it. If the present account agrees with all of that, what does it add to Recanati? The answer marks the one place the two accounts genuinely part. Recanati's file is a vehicle --- a mental particular that carries information and is individuated by the relation it exploits; the relation is named as epistemically rewarding and is left there, as a primitive the file stands on. The present account does not posit a vehicle. What anchors a name is not a file but a coordination --- an \(x^o\)-grip on an encountered particular --- that is an episode in a practice rather than a standing mental object, and that the account describes rather than names: it can misfire, it is corrected, it is the sort of thing one does correctly or incorrectly. The difference is between a posited container whose individuating relation is a primitive and a described operation whose normativity is on the surface. Where Recanati locates singular reference in the file and treats the relation as a given, the present account locates it in the relation and treats the file's work --- the holding-fixed --- as the thing to be explained. That is not a relabelling of mental files; it is a different answer to what the file was introduced to do.

\subsection{The Sophisticated-Descriptivism Objection}\label{the-sophisticated-descriptivism-objection}

The objection. This is the sharper charge. The account concedes that contemporary speakers have no R-grounding with Aristotle, that they operate through A-grounding, and that what they cognitively possess is a bundle of descriptions \(\mathcal{D}_1, \mathcal{D}_2, \mathcal{D}_3\) coordinated with the name. But this is what descriptivism always said: that competent use of a name consists in associating descriptions with it. Dressing the bundle as a conjunction of revisable attributions \(\mathcal{R}(\mathcal{D}_i, N^c)\), and insisting that the descriptions are ``attributed to'' rather than ``constitutive of'' the referent, is a distinction without a difference. If you ask what, in the contemporary speaker's head, fixes which individual she is talking about, the only available answer on this account is: the descriptions. So reference is carried by the descriptions, and the anti-descriptivist protestation is empty.

Reply. The objection mislocates where reference-fixing happens, and this mislocation is exactly what the account is built to expose. The question ``What, in the later speaker's mind, fixes the referent?'' falsely presupposes that reference is fixed by the current user. On the present account, it is not. Reference was fixed once, at baptism, by a demonstrative coordination that downstream speakers neither performed nor can perform. What is in her head is not a reference-fixing device but an inheritance: the name-concept \(N^c\), acquired through the chain, already carrying the reference that the baptismal \(\mathcal{R}(h^o, N^c)\) established. The descriptions she attaches to \(N^c\) are how she cognitively inhabits an inheritance she did not author. They are not a mechanism for selecting a referent; they are what she has come to believe about a referent already selected for her.

This is not a verbal difference, because the two pictures come apart on the data. Take the case the objection cannot absorb: the speaker whose every description is false. Suppose a student associates with ``Aristotle'' only ``the author of the \emph{Nicomachean Ethics}'' and ``Plato's most famous pupil,'' and suppose --- counterfactually --- that scholarship overturns both, showing the \emph{Ethics} was compiled by students and that the famous pupil was someone else. On descriptivism, the student's ``Aristotle'' now refers to whoever satisfies the surviving descriptions, or to no one. On the present account, the student still refers to Aristotle, because reference rode on \(N^c\) the whole time and \(N^c\) traces through the chain to the baptismal grounding, not through the student's beliefs. The descriptions were attributions; attributions can be wholly false without disturbing reference, exactly as a person's reputation can be entirely mistaken without making the reputation about someone else. Descriptivism cannot say this. The account predicts it. That is the difference the objection says does not exist.

The deeper point is structural. Descriptivism makes the descriptions do two jobs at once --- fix the referent and express the speaker's understanding of it --- and the anti-descriptivist arguments (Kripke's Feynman, Gödel) are arguments that the descriptions cannot do the first job. The present account agrees they cannot, and so it removes the first job from the descriptions entirely and assigns it to the baptismal grounding transmitted as \(N^c\). The descriptions keep only the second job. A view that gives descriptions one of descriptivism's two jobs is not descriptivism, any more than a view on which a portrait resembles its sitter without determining who the sitter is collapses into the view that the sitter just is whoever the portrait resembles.

One residue should be granted rather than papered over. Reference rides on \(N^c\) through the chain --- but a later speaker's grip on which chain her use of ``Aristotle'' belongs to is itself secured partly through testimonial and descriptive material (``the Aristotle my teacher meant,'' ``the philosopher, not the shipping magnate''). The descriptivist can relocate the worry here: not at reference-fixing, but at chain-individuation. This is a real cost, but a smaller one. Disambiguating which chain one is on is an epistemic task that draws on description; it does not make description constitutive of the reference the chain carries, any more than needing a catalogue number to find a painting makes the number constitutive of which painting it is. The worry moves from the center of the account to its edge, and changes from ``reference is descriptive'' to ``chain-identification is sometimes descriptively mediated,'' which the account can grant without giving up anti-descriptivism about reference itself.

\subsection{The Unargued-Categorial-Feature Objection}\label{the-unargued-categorial-feature-objection}

The objection. The account's treatment of absence rests on a claim stated in the introduction and never defended: that proper names, as a category of linguistic item, are treated by competent speakers as tracking unique individuals who actually existed and were originally baptized. This categorial feature is what carries reference past the point where R-grounding lapses. But the feature is exactly what needed explaining. To say that names persist in reference because they belong to a category whose members persist in reference is to name the explanandum and call it an explanans. The wheel turns nothing.

Reply. The charge is fair as stated, and meeting it requires saying what the categorial feature is and why positing it is not circular. The feature is not a brute semantic stipulation that names refer rigidly --- that would indeed be the explanandum relabelled. It is a fact about a practice: the community treats name-tokens as governed by a norm of co-reference with their baptismal origin, and enforces that norm through correction. When a child uses ``Aristotle'' of the wrong figure, she is corrected, and the correction does not cite descriptive failure (``Aristotle wrote no such thing'') but identity failure (``that's not who `Aristotle' picks out''). The categorial feature is the standing of this norm in the practice, and it is no more circular to appeal to it than it is circular to explain the rigidity of a measurement standard by appeal to the institution that maintains it.

What keeps the appeal non-circular is that the norm has an independent characterization and an independent origin. Its characterization is behavioral and observable: it is the pattern of correction, deference, and chain-preserving transmission that Kripke described from the outside and that the constructive chapter redescribes in cognitive terms. Its origin is the baptismal grounding: the norm of co-reference-with-origin can attach to a name precisely because there was an origin --- a real coordination with a real individual --- for the norm to point back to. A category of terms with no such origins (genuinely empty names aside, flagged as an open case in the conclusion) could not sustain the norm, which is why descriptions, which have no baptismal origin, do not behave as names do. So the categorial feature is not posited to explain rigidity and then used to explain rigidity. It is identified independently --- as a correction-practice with a baptismal anchor --- and then shown to deliver persistence across absence. The explanation runs from practice-plus-origin to persistence, not from persistence to itself.

I do not claim this dissolves every form of the worry. An objector may press: why does this practice exist, why do human communities maintain origin-tracking norms for names but not for descriptions? That is a genuine question, and it is anthropological and developmental as much as philosophical; the account locates the answer in the usefulness of carrying a fixed grip past the encounter but does not pretend to a full genealogy. What the reply establishes is narrower and sufficient for the thesis: the categorial feature is a describable practice with a describable anchor, not a synonym for the conclusion.

\subsection{\texorpdfstring{The ``\(x^o\) Is Just Perception'' Objection}{The ``x\^{}o Is Just Perception'' Objection}}\label{the-xo-is-just-perception-objection}

The objection. The whole construction depends on \(x^o\) being a mode of language --- a word-use --- rather than a relabelled perceptual or causal relation between a speaker and an object. If \(x^o\) is really just perception (or perceptual attention, or a causal-informational link), then the account is a causal theory of reference in disguise, and the coordination \(\mathcal{R}(bird^o, E\{BIRD\})\) is just the old causal relation with a new symbol. The claim that \(x^o\) ``is still language, not perception'' is asserted in the constructive chapter but, the objection runs, cannot be argued, because there is nothing for \(x^o\) to be between perception and inferential role.

Reply. \(x^o\) is distinguished from perception by a double dissociation, and the dissociation is what shows it is a third thing rather than either pole. On one side, the experience can occur without the word: \(E\{X\}\) --- the perceptual-motor-affective activation --- happens in pre-linguistic infants and in animals with no word at all. Perception does not require \(x^o\). On the other side, the word can occur, in object-mode, and misfire on its object: a speaker says ``that bird'' while pointing, and coordinates with the wrong thing in the visual field, or with a leaf she has mistaken for a bird. The word's object-mode use can be wrong about what it coordinates with, which a perceptual state cannot be in the same way --- the perception is whatever it is; the word applied to it is what bears correctness. That a use can be incorrect is the mark of the normative, and the normative is the mark of language. So \(x^o\) is not perception (perception runs without it) and not free of the world (it can misfire against the world), and the space between those two facts is exactly the space \(x^o\) occupies: a public, correctable word-use that requires coordination with what is encountered but is not identical to the encounter.

This also answers the causal-theory charge. A causal-informational link between word and object is not the sort of thing that can be mistaken while remaining the same link; it either obtains or does not. \(x^o\) can be exercised incorrectly --- wrong instance, wrong sortal --- and be recognized as incorrect by the speaker and the community, and corrected. Correctability is not a feature of causal relations; it is a feature of moves in a practice. The coordination \(\mathcal{R}\) is therefore not the causal relation relabelled, because it carries a normative status the causal relation lacks. What \(x^o\) shares with the causal theory is only the insistence that reference reaches outside the linguistic network; what it adds, and what makes it a linguistic mode rather than a causal hook, is that the reaching-out is something one can do correctly or incorrectly.

\subsection{Two Smaller Objections}\label{two-smaller-objections}

Demonstrative misfire and the rigidity claim. If \(x^o\) can misfire, how can local rigidity be rigid? A designator that might be coordinating with the wrong object hardly seems to ``hold its referent fixed.'' The reply is that rigidity and infallibility are the two properties the dissertation has been at pains to separate since the chapter on Russell. What a misfire puts at risk is the grounding --- which particular, if any, the coordination has caught. What rigidity governs is the behavior, under modal variation, of whatever the coordination has in fact caught. Given that the use coordinates with this bird, every counterfactual built on it is about this bird; that is the rigidity, and it is untouched by the fact that the speaker could have been coordinating with a different bird, or none. The misfire is a mistake about the target; the rigidity is the modal profile of the target once fixed. These are orthogonal, and conflating them is the very error that drove Russell to immediate sense-data.

The status of surrogate objects. The constructive account grants that a blind historian refers to Aristotle with no surrogate \(\sigma\), which seems to make \(\sigma\) idle. The reply is that \(\sigma\) is not part of the reference-fixing mechanism and was never claimed to be; it is one of three transmission channels, and the weakest. Testimony \(\tau\) and the linguistic pattern \(\langle N \rangle\) carry reference; surrogates only enrich the speaker's phenomenological relation to the bearer. The blind historian shows that \(\sigma\) is dispensable, not that it is doing covert work. An account on which surrogates were load-bearing would be embarrassed by the blind historian; this account predicts her, because it locates reference in the chain and the norm, not in any residual perceptual contact.

\subsection{What the Objections Leave Standing}\label{what-the-objections-leave-standing}

The four principal objections share a structure: each assumes that the account must locate reference-fixing in the contemporary individual speaker --- in her distinction-drawing, her descriptions, her perceptual link, or a brute semantic feature she deploys. The thesis's recurring reply is that reference-fixing happened once, at a baptismal grounding, and that everything the contemporary speaker does is inheritance and inhabitation of a grip established for her. The restatement objection misses the direction of derivation; the descriptivism objection misplaces reference in the speaker's beliefs; the categorial-feature objection mistakes a describable practice for a relabelled conclusion; the perception objection collapses a correctable word-use into an uncorrectable link. Once reference-fixing is put back at the origin, where the account places it, each objection loses its grip. None of this shows the account is true. It shows the account is not any of the four things it is most likely to be mistaken for.

\section{Conclusion}\label{conclusion}

\subsection{What Has Been Argued}\label{what-has-been-argued}

The dissertation began from a gap in Kripke. The semantic results of \emph{Naming and Necessity} tell us what rigid designators do --- hold their referent fixed across every world in which it exists --- and Kripke argues for those results on grounds that owe nothing to the necessity-of-identity formula. What the results do not tell us is what cognitive operation a speaker performs such that her words behave as the semantics describes. Kripke saw that reference to a particular cannot be secured by formal operations alone, which is why he appealed to baptism and the causal chain; but he offered that appeal as a better picture, not a theory, and left the operation it pictures undescribed. The dissertation set out to describe it.

The argument was built in sequence. It opened by setting out the field --- direct-reference theory and Donnellan's two contributions --- and by separating its own question, the single speaker's grip on a particular, from the inter-speaker coordination question it set aside. It then stated Kripke's rigid-designation thesis in his own terms, and located the cognitive gap not in the thesis but at a specific point in the necessity-of-identity argument: the word ``object,'' in the formula's natural-language gloss, spans variable, encountered particular, and named individual, and the formula's validity holds for the first while its philosophical payload concerns the third. The move between the levels is where a cognitive operation is presupposed and not supplied. The survey of existing attempts showed that the operation had been approached and missed: Russell located it --- direct, non-descriptive coordination --- and then confined it to incorrigible sense-data by tying directness to infallibility, while Evans, Recanati, and Dickie each captured part of the burden without describing the operation itself.

The constructive chapters named and developed it. A word in object-mode coordinates with an encountered particular; that coordination grounds the use, and rigidity is what the grip amounts to under modal variation. Baptism stabilizes such a coordination into a public name --- not transferring a ready-made property, but giving a singular anchoring a durable vehicle. Testimonial transmission carries the name across speakers and generations who never had the encounter, and contemporary speakers engage absent bearers through revisable attributions that ride on a reference they do not constitute. The objections the account most invites --- that it relabels an old distinction, that it is descriptivism in disguise, that it helps itself to an unargued feature of names, that its object-mode is only perception renamed --- were answered by locating reference-fixing at the origin rather than in the contemporary speaker.

The single claim the dissertation has tried to establish is that Kripkean rigidity is not primitive. It is local rigidity made portable: the holding-fixed that a demonstrative coordination has within an encounter, carried past the encounter by a social practice that did not create it but only changed its locus. This is a metasemantic claim, not a rival semantics. It leaves Kripke's model untouched and explains how a speaker comes to be positioned such that the model correctly describes her practice.

\subsection{What Remains Open}\label{what-remains-open}

The thesis was deliberately narrow, and its scope leaves several questions standing. I mark them, since an account is more useful with its edges visible.

\emph{General-term rigidity.} The object-mode / concept-mode distinction applies to common nouns as well as names, and the account therefore has the materials to ask whether and how ``water'' or ``gold'' are rigid. But general-term rigidity has its own large literature and its own complications --- the role of scientific theorizing in fixing kind-reference, the metaphysics of natural kinds --- on which the dissertation took no stand. Whether local rigidity scales to kinds the way it scales to individuals is the most natural extension of the account.

\emph{Trace-baptism and empty names.} The account of baptism took as its paradigm the case in which the baptizer is in object-mode coordination with the bearer, and marked the contrasting structure where a name is introduced through a description of an unencountered cause --- ``Neptune,'' ``Jack the Ripper'' --- as a coordination of a different argument type. Whether such a baptism anchors its name as ostensive baptism does, or only through the naming norm, was left open, and with it the treatment of empty and fictional names, which would differ in mechanism: a trace-baptism whose causal inference misfired, as with ``Vulcan,'' differs from a baptism made within a pretense.

\emph{The unity of holding-fixed.} The derivation claim was scoped to names of encountered particulars, granting that holding an object fixed for de re predication can also be secured within the linguistic-formal system, as with the number read de re or with mathematical reference, where there is no encounter to ground a coordination. Whether these are exercises of a single capacity --- the tracking capacity redeployed on targets secured within language --- or two different operations was left as a question the dissertation is not positioned to settle, since it turns partly on developmental evidence about whether descriptive fixation is parasitic on prior tracking.

\emph{Mathematical reference.} Mathematical objects are rigid, yet there is no perceptual encounter to ground an object-mode coordination with a number. The account treats such fixation as secured within the formal system rather than by reaching outside it; whether that is the whole story is left open.

\emph{The diagnostic application to Russell, Quine, and Lewis.} The suggestion that a shared over-reliance on formal operations runs through Russell's sense-data, Quine's bound variables, and Lewis's plurality of worlds, and that the grounding / non-grounding distinction diagnoses it, was used for Russell and reserved, for the others, to separate work that would test how far the distinction generalizes.

\emph{Inter-speaker coordination.} The dissertation isolated the single speaker's anchoring of a word to a particular and set aside the distinct question of how a community of speakers stays coordinated on a shared referent --- the phenomenon the relationist tradition takes as central. The account's claim was one of order: that coordination operates on a name already anchored, not that coordination is where reference is fixed. Defending that order against a developed relationism, and giving the inter-speaker practice its own treatment, is a connected project reserved for separate work.

\emph{Grounding as the primitive.} The account was organized around a directional claim it used but did not argue: that the coordination of a word with a particular comes first, and that rigidity is the modal manifestation of that coordination. Taking this further --- analyzing what grounding is, asking whether rigidity is a consequence of grounding rather than a notion standing alongside it, and characterizing the coordination operator accordingly as a grounding relation --- is a project of its own, requiring a conceptual analysis the present argument neither needed nor could carry. It is the natural development of the account beyond the dissertation.

\subsection{The Shape of the Result}\label{the-shape-of-the-result}

What the dissertation offers is small and specific: not a theory of reference, but the closing of one gap in an existing theory. Kripke described a practice and declined to analyze the operation it rests on. The operation, I have argued, is one we can watch in its simplest form whenever a speaker points at a bird and says something about how it might have been --- the holding-fixed of a particular across everything imagined about it. Names are that holding-fixed, made to survive the absence of the thing held. The cognitive precondition Kripke presupposed is the same one Russell found and abandoned; what the account does is take it back from the place Russell left it and follow it to where Kripke needed it.

\section{Bibliography}\label{bibliography}

\subsection{Primary Literature}\label{primary-literature}

\emph{(Texts the dissertation directly analyzes or argues against)}

\textbf{Donnellan, K.} (1966). ``Reference and Definite Descriptions.''~\emph{Philosophical Review}, 75(3), 281--304.

\textbf{Donnellan, K.} (1970). ``Proper Names and Identifying Descriptions.''~\emph{Synthese}, 21(3--4), 335--358.

\textbf{Kaplan, D.} (1968). ``Quantifying In.'' \emph{Synthese}, 19(1--2), 178--214.

\textbf{Kaplan, D.} (1989). ``Demonstratives.'' In J. Almog, J. Perry \& H. Wettstein (eds.), \emph{Themes from Kaplan}. Oxford: Oxford University Press, pp.~481--563.

\textbf{Kripke, S. A.} (1971). ``Identity and Necessity.'' In M. K. Munitz (ed.), \emph{Identity and Individuation}. New York: New York University Press, pp.~135--164.

\textbf{Kripke, S. A.} (1980). \emph{Naming and Necessity}. Cambridge, MA: Harvard University Press.

\textbf{Quine, W. V. O.} (1953). ``Reference and Modality.'' In \emph{From a Logical Point of View}. Cambridge, MA: Harvard University Press, pp.~139--159.

\textbf{Quine, W. V. O.}~(1953). ``Three Grades of Modal Involvement.''~\emph{Proceedings of the XIth International Congress of Philosophy}, 14, 65--81. {[}Repr. in~\emph{The Ways of Paradox}. New York: Random House, 1966, pp.~156--174.{]}

\textbf{Russell, B.} (1905). ``On Denoting.'' \emph{Mind}, 14(56), 479--493.

\textbf{Russell, B.}~(1918/1956). ``The Philosophy of Logical Atomism.'' In R. C. Marsh (ed.),~\emph{Logic and Knowledge}. London: Allen \& Unwin, pp.~177--281.

\textbf{Wittgenstein, L.}~(1922).~\emph{Tractatus Logico-Philosophicus}. Trans. C. K. Ogden, with assistance from F. P. Ramsey. London: Kegan Paul, Trench, Trubner \& Co.

\textbf{Wittgenstein, L.} (1953). \emph{Philosophical Investigations}. Trans. G. E. M. Anscombe. Oxford: Blackwell.

\subsection{Secondary Literature}\label{secondary-literature}

\textbf{Ayer, A. J.}~(1936).~\emph{Language, Truth and Logic}. London: Victor Gollancz.

\textbf{Carnap, R.}~(1947).~\emph{Meaning and Necessity: A Study in Semantics and Modal Logic}. Chicago: University of Chicago Press.

\textbf{Chalmers, D. J.}~(2006). ``Two-Dimensional Semantics.'' In E. Lepore \& B. C. Smith (eds.),~\emph{The Oxford Handbook of Philosophy of Language}. Oxford: Oxford University Press, pp.~575--606.

\textbf{Davies, M. \& Humberstone, L.}~(1980). ``Two Notions of Necessity.''~\emph{Philosophical Studies}, 38(1), 1--30.

\textbf{Dickie, I.} (2015).~\emph{Fixing Reference}. Oxford: Oxford University Press.

\textbf{Evans, G.} (1973). ``The Causal Theory of Names.''~\emph{Proceedings of the Aristotelian Society}, Supplementary Vol. 47, 187--208.

\textbf{Evans, G.} (1982). \emph{The Varieties of Reference}. Oxford: Oxford University Press.

\textbf{Fine, K.} (2007).~\emph{Semantic Relationism}. Oxford: Blackwell.

\textbf{Frege, G.}~(1891). ``Function and Concept.'' In P. Geach \& M. Black (eds.),~\emph{Translations from the Philosophical Writings of Gottlob Frege}. Oxford: Blackwell, 1952, pp.~21--41.

\textbf{Frege, G.}~(1892). ``On Sense and Reference.'' In P. Geach \& M. Black (eds.),~\emph{Translations from the Philosophical Writings of Gottlob Frege}. Oxford: Blackwell, 1952, pp.~56--78.

\textbf{LaPorte, J.} (2022). ``Rigid Designators.'' In E. N. Zalta (ed.), \emph{The Stanford Encyclopedia of Philosophy} (Spring 2022 ed.). \url{https://plato.stanford.edu/entries/rigid-designators/}

\textbf{Lewis, D.}~(1986).~\emph{On the Plurality of Worlds}. Oxford: Blackwell.

\textbf{Putnam, H.} (1975). ``The Meaning of `Meaning'.'' In \emph{Mind, Language and Reality: Philosophical Papers, Vol. 2}. Cambridge: Cambridge University Press, pp.~215--271. {[}Also \emph{Minnesota Studies in the Philosophy of Science}, 7, 131--193.{]}

\textbf{Recanati, F.} (2012). \emph{Mental Files}. Oxford: Oxford University Press.

\textbf{Salmon, N.} (1981). \emph{Reference and Essence}. Princeton: Princeton University Press.

\textbf{Smullyan, A. F.} (1948). ``Modality and Description.'' \emph{Journal of Symbolic Logic}, 13(1), 31--37.

\textbf{Soames, S.} (2002). \emph{Beyond Rigidity: The Unfinished Semantic Agenda of} Naming and Necessity. New York: Oxford University Press.

\subsection{Author's Own Work}\label{authors-own-work}

\textbf{Cojocariu, F.} (2026). ``Pattern-Recognition Unity: A Framework Specification.'' Preprint. \url{https://dx.doi.org/10.2139/ssrn.6285878}

\vfill
\noindent\rule{\textwidth}{0.4pt}\\[2pt]
{\tiny florin.cojocariu@s.unibuc.ro, \today}

\end{document}